
\documentclass[12pt, a4paper]{article}

% ==========================================
% 1. COMPILER, LANGUAGE & FONT SETUP
% ==========================================
\usepackage{fontspec}
\usepackage{polyglossia}
\usepackage[protrusion=false]{microtype}
\setmainlanguage{bengali}
\setotherlanguage{english}
\newfontfamily\bengalifont[Script=Bengali, AutoFakeBold=2.5, AutoFakeSlant=0.2]{Kalpurush}
\usepackage{enumitem}
\setlist[itemize,1]{label=\textendash}

% ==========================================
% 2. MATHEMATICS, UNITS & FORMATTING
% ==========================================
\usepackage{amsmath, amssymb, siunitx}
\usepackage[margin=1in]{geometry}
\usepackage{setspace}
\usepackage{enumitem}
\usepackage{microtype} % For better typography
\onehalfspacing % Better line spacing for readability

% ==========================================
% 3. COLORS & BEAUTIFICATION PACKAGES
% ==========================================
\usepackage[dvipsnames, table]{xcolor}
\usepackage[skins, breakable, theorems]{tcolorbox}
\usepackage{tikz}
\usetikzlibrary{shadows, arrows.meta, positioning, decorations.pathmorphing, calc}

% Define custom beautiful colors
\definecolor{primary}{HTML}{0056b3}
\definecolor{secondary}{HTML}{e7f1ff}
\definecolor{success}{HTML}{198754}
\definecolor{successbg}{HTML}{d1e7dd}
\definecolor{accent}{HTML}{fd7e14}
\definecolor{darkgray}{HTML}{343a40}

% Custom Tcolorbox Styles
\tcbset{
    questionbox/.style={
        enhanced, colback=secondary, colframe=primary, arc=2mm, boxrule=1pt,
        fonttitle=\bfseries\Large, title=#1,
        drop fuzzy shadow=black!10,
        attach boxed title to top left={xshift=5mm, yshift=-3mm},
        boxed title style={colback=primary, arc=1mm}
    },
    answerbox/.style={
        enhanced, colback=successbg, colframe=success, arc=1.5mm, boxrule=1pt,
        left=2mm, right=2mm, top=2mm, bottom=2mm,
        drop fuzzy shadow=black!10
    },
    solutionbox/.style={
        enhanced, colback=white, colframe=darkgray, arc=2mm, boxrule=1pt,
        fonttitle=\bfseries\large, title=#1, breakable,
        coltitle=white, colbacktitle=darkgray,
        drop fuzzy shadow=black!10,
        attach boxed title to top center={yshift=-3mm},
        boxed title style={arc=1mm}
    }
}

\begin{document}

% ------------------------------------------
% QUESTION SECTION
% ------------------------------------------
\begin{tcolorbox}[questionbox={প্রশ্ন (Question) 1}]
    \noindent
    $x^2 + y^2 - 2x - 4y - 20 = 0$ এবং $x^2 + y^2 - 14x - 14y + 94 = 0$ বৃত্তদ্বয়ের সাধারণ স্পর্শকের সংখ্যা কত এবং তাদের বাহ্যিক সমরূপতা কেন্দ্রের (External Center of Similitude) স্থানাঙ্ক কোনটি?
    
    \vspace{0.8em}
    \begin{english}
    \noindent\textit{\color{darkgray}
    What is the number of common tangents to the circles $x^2 + y^2 - 2x - 4y - 20 = 0$ and $x^2 + y^2 - 14x - 14y + 94 = 0$, and what are the coordinates of their external center of similitude?
    }
    \end{english}
\end{tcolorbox}

% ------------------------------------------
% OPTIONS SECTION
% ------------------------------------------
\begin{itemize}[label={}, leftmargin=1cm, itemsep=0.5em]
    \item[\textbf{A)}] $4$ টি সাধারণ স্পর্শক, বাহ্যিক সমরূপতা কেন্দ্র $(11, \frac{31}{3})$
    \item[\textbf{B)}] $2$ টি সাধারণ স্পর্শক, বাহ্যিক সমরূপতা কেন্দ্র $(11, 10)$
    \item[\textbf{C)}] $4$ টি সাধারণ স্পর্শক, বাহ্যিক সমরূপতা কেন্দ্র $(11, 10)$
    \item[\textbf{D)}] $3$ টি সাধারণ স্পর্শক, বাহ্যিক সমরূপতা কেন্দ্র $(11, \frac{31}{3})$
\end{itemize}

% ------------------------------------------
% ANSWER HIGHLIGHT
% ------------------------------------------
\vspace{0.5em}
\begin{tcolorbox}[answerbox]
    \textbf{\color{success} উত্তর:} A) $4$ টি সাধারণ স্পর্শক, বাহ্যিক সমরূপতা কেন্দ্র $(11, \frac{31}{3})$
\end{tcolorbox}
\vspace{1em}

% ------------------------------------------
% SOLUTION SECTION
% ------------------------------------------
\begin{tcolorbox}[solutionbox={সমাধান (Solution)}]
    \noindent
    প্রথম বৃত্ত: $S_1 = x^2 + y^2 - 2x - 4y - 20 = 0$ এর কেন্দ্র $C_1(1, 2)$ এবং ব্যাসার্ধ 
    \[ r_1 = \sqrt{1^2 + 2^2 - (-20)} = \sqrt{25} = 5 \]
    দ্বিতীয় বৃত্ত: $S_2 = x^2 + y^2 - 14x - 14y + 94 = 0$ এর কেন্দ্র $C_2(7, 7)$ এবং ব্যাসার্ধ 
    \[ r_2 = \sqrt{(-7)^2 + (-7)^2 - 94} = \sqrt{49 + 49 - 94} = \sqrt{4} = 2 \]
    বৃত্তদ্বয়ের কেন্দ্রদ্বয়ের মধ্যবর্তী দূরত্ব $d = C_1C_2$:
    \[ d = \sqrt{(7 - 1)^2 + (7 - 2)^2} = \sqrt{6^2 + 5^2} = \sqrt{36 + 25} = \sqrt{61} \approx 7.81 \]
    যেহেতু $d = \sqrt{61} > r_1 + r_2 = 7$, বৃত্তদ্বয় সম্পূর্ণ পৃথক এবং একে অপরকে ছেদ বা স্পর্শ করে না। অতএব, এদের সাধারণ স্পর্শকের সংখ্যা $4$ টি।
    
    \vspace{0.5em}
    আমরা জানি, বাহ্যিক সমরূপতা কেন্দ্র $E(x_e, y_e)$ বৃত্তদ্বয়ের কেন্দ্রদ্বয়কে $r_1 : r_2 = 5 : 2$ অনুপাতে বহিঃস্থভাবে বিভক্ত করে:
    \begin{align*}
        x_e &= \frac{r_1 x_2 - r_2 x_1}{r_1 - r_2} = \frac{5(7) - 2(1)}{5 - 2} = \frac{35 - 2}{3} = \frac{33}{3} = 11 \\
        y_e &= \frac{r_1 y_2 - r_2 y_1}{r_1 - r_2} = \frac{5(7) - 2(2)}{5 - 2} = \frac{35 - 4}{3} = \frac{31}{3}
    \end{align*}
    অতএব, স্পর্শকের সংখ্যা $4$ এবং বাহ্যিক সমরূপতা কেন্দ্রের স্থানাঙ্ক $\left(11, \frac{31}{3}\right)$।
\end{tcolorbox}

% ------------------------------------------
% QUESTION SECTION
% ------------------------------------------
\begin{tcolorbox}[questionbox={প্রশ্ন (Question) 2}]
    \noindent
    একটি বৃত্ত $x^2 + y^2 - 4x - 6y - 12 = 0$ বৃত্তের সাথে সমকেন্দ্রিক এবং এটি $3x + 4y - 5 = 0$ সরলরেখা থেকে এমন একটি জ্যা কর্তন করে যার দৈর্ঘ্য $8$ একক। বৃত্তটির সমীকরণ কোনটি?
    
    \vspace{0.8em}
    \begin{english}
    \noindent\textit{\color{darkgray}
    A circle is concentric with the circle $x^2 + y^2 - 4x - 6y - 12 = 0$ and intercepts a chord of length $8$ units on the line $3x + 4y - 5 = 0$. Which of the following is the equation of the circle?
    }
    \end{english}
\end{tcolorbox}

% ------------------------------------------
% OPTIONS SECTION
% ------------------------------------------
\begin{itemize}[label={}, leftmargin=1cm, itemsep=0.5em]
    \item[\textbf{A)}] $25(x^2 + y^2) - 100x - 150y - 244 = 0$
    \item[\textbf{B)}] $25(x^2 + y^2) - 100x - 150y - 250 = 0$
    \item[\textbf{C)}] $5(x^2 + y^2) - 20x - 30y - 48 = 0$
    \item[\textbf{D)}] $25(x^2 + y^2) - 50x - 75y - 122 = 0$
\end{itemize}

% ------------------------------------------
% ANSWER HIGHLIGHT
% ------------------------------------------
\vspace{0.5em}
\begin{tcolorbox}[answerbox]
    \textbf{\color{success} উত্তর:} A) $25(x^2 + y^2) - 100x - 150y - 244 = 0$
\end{tcolorbox}
\vspace{1em}

% ------------------------------------------
% SOLUTION SECTION
% ------------------------------------------
\begin{tcolorbox}[solutionbox={সমাধান (Solution)}]
    \noindent
    যেহেতু বৃত্তটি $x^2 + y^2 - 4x - 6y - 12 = 0$ বৃত্তের সাথে সমকেন্দ্রিক, তাই বৃত্তটির কেন্দ্র হবে প্রদত্ত বৃত্তের কেন্দ্রের সমান। কেন্দ্র $C(2, 3)$। ধরি, নির্ণেয় বৃত্তের ব্যাসার্ধ $R$। 
    
    \vspace{0.5em}
    কেন্দ্র $C(2, 3)$ হতে সরলরেখা $3x + 4y - 5 = 0$ এর লম্ব দূরত্ব $d$:
    \[ d = \frac{|3(2) + 4(3) - 5|}{\sqrt{3^2 + 4^2}} = \frac{|6 + 12 - 5|}{5} = \frac{13}{5} = 2.6 \]
    বৃত্তটি দ্বারা রেখাটি হতে কর্তিত জ্যা-এর দৈর্ঘ্য $8$ একক হলে, জ্যামিতিক সম্পর্ক হতে পাই:
    \[ 2\sqrt{R^2 - d^2} = 8 \implies R^2 - d^2 = 16 \]
    \[ R^2 = 16 + d^2 = 16 + \left(\frac{13}{5}\right)^2 = 16 + \frac{169}{25} = \frac{400 + 169}{25} = \frac{569}{25} \]
    বৃত্তটির সমীকরণ:
    \[ (x - 2)^2 + (y - 3)^2 = R^2 \]
    \[ (x - 2)^2 + (y - 3)^2 = \frac{569}{25} \]
    \[ x^2 - 4x + 4 + y^2 - 6y + 9 = \frac{569}{25} \]
    \[ x^2 + y^2 - 4x - 6y + 13 = \frac{569}{25} \]
    উভয়পক্ষকে $25$ দ্বারা গুণ করে পাই:
    \[ 25(x^2 + y^2 - 4x - 6y + 13) = 569 \]
    \[ 25(x^2 + y^2) - 100x - 150y + 325 - 569 = 0 \]
    \[ 25(x^2 + y^2) - 100x - 150y - 244 = 0 \]
    অতএব, বৃত্তটির সমীকরণ $25(x^2 + y^2) - 100x - 150y - 244 = 0$।
\end{tcolorbox}

% ------------------------------------------
% QUESTION SECTION
% ------------------------------------------
\begin{tcolorbox}[questionbox={প্রশ্ন (Question) 3}]
    \noindent
    $x^2 + y^2 - 2x - 4y + 1 = 0$ বৃত্তের উপর অবস্থিত যেকোনো বিন্দু $P(x, y)$ এর জন্য $3x - 4y$ রাশিটির সর্বোচ্চ এবং সর্বনিম্ন মানের গুণফল কত?
    
    \vspace{0.8em}
    \begin{english}
    \noindent\textit{\color{darkgray}
    For any point $P(x, y)$ on the circle $x^2 + y^2 - 2x - 4y + 1 = 0$, what is the product of the maximum and minimum values of the expression $3x - 4y$?
    }
    \end{english}
\end{tcolorbox}

% ------------------------------------------
% OPTIONS SECTION
% ------------------------------------------
\begin{itemize}[label={}, leftmargin=1cm, itemsep=0.5em]
    \item[\textbf{A)}] $-75$
    \item[\textbf{B)}] $-50$
    \item[\textbf{C)}] $75$
    \item[\textbf{D)}] $-100$
\end{itemize}

% ------------------------------------------
% ANSWER HIGHLIGHT
% ------------------------------------------
\vspace{0.5em}
\begin{tcolorbox}[answerbox]
    \textbf{\color{success} উত্তর:} A) $-75$
\end{tcolorbox}
\vspace{1em}

% ------------------------------------------
% SOLUTION SECTION
% ------------------------------------------
\begin{tcolorbox}[solutionbox={সমাধান (Solution)}]
    \noindent
    প্রদত্ত বৃত্তের সমীকরণ: $x^2 + y^2 - 2x - 4y + 1 = 0$
    বৃত্তের কেন্দ্র $C(1, 2)$ এবং ব্যাসার্ধ $r = \sqrt{1^2 + 2^2 - 1} = 2$।
    
    \vspace{0.5em}
    ধরি, $3x - 4y = k \implies 3x - 4y - k = 0$। যেহেতু $P(x, y)$ বিন্দুটি বৃত্তের উপর অবস্থিত, তাই রেখাটি বৃত্তকে ছেদ বা স্পর্শ করবে। স্পর্শ বা ছেদের শর্ত হলো, কেন্দ্র হতে রেখাটির লম্ব দূরত্ব ব্যাসার্ধের চেয়ে ছোট বা সমান হবে:
    \[ \frac{|3(1) - 4(2) - k|}{\sqrt{3^2 + 4^2}} \le 2 \]
    \[ \frac{|3 - 8 - k|}{5} \le 2 \]
    \[ |-5 - k| \le 10 \implies |k + 5| \le 10 \]
    \[ -10 \le k + 5 \le 10 \implies -15 \le k \le 5 \]
    সুতরাং, রাশিটির সর্বোচ্চ মান $k_{\max} = 5$ এবং সর্বনিম্ন মান $k_{\min} = -15$। 
    সর্বোচ্চ ও সর্বনিম্ন মানের গুণফল:
    \[ k_{\max} \times k_{\min} = 5 \times (-15) = -75 \]
    অতএব, গুণফলটি $-75$।
\end{tcolorbox}

% ------------------------------------------
% QUESTION SECTION
% ------------------------------------------
\begin{tcolorbox}[questionbox={প্রশ্ন (Question) 4}]
    \noindent
    পোলার সমীকরণ $r = 6\cos\theta + 8\sin\theta$ বৃত্তটি দ্বারা $\theta = \frac{\pi}{4}$ রেখার উপর কর্তিত জ্যা-এর দৈর্ঘ্য কত একক?
    
    \vspace{0.8em}
    \begin{english}
    \noindent\textit{\color{darkgray}
    What is the length of the chord in units intercepted by the circle $r = 6\cos\theta + 8\sin\theta$ on the line $\theta = \frac{\pi}{4}$?
    }
    \end{english}
\end{tcolorbox}

% ------------------------------------------
% OPTIONS SECTION
% ------------------------------------------
\begin{itemize}[label={}, leftmargin=1cm, itemsep=0.5em]
    \item[\textbf{A)}] $7\sqrt{2}$ একক
    \item[\textbf{B)}] $5\sqrt{2}$ একক
    \item[\textbf{C)}] $6\sqrt{2}$ একক
    \item[\textbf{D)}] $10$ একক
\end{itemize}

% ------------------------------------------
% ANSWER HIGHLIGHT
% ------------------------------------------
\vspace{0.5em}
\begin{tcolorbox}[answerbox]
    \textbf{\color{success} উত্তর:} A) $7\sqrt{2}$ একক
\end{tcolorbox}
\vspace{1em}

% ------------------------------------------
% SOLUTION SECTION
% ------------------------------------------
\begin{tcolorbox}[solutionbox={সমাধান (Solution)}]
    \noindent
    বৃত্তের পোলার সমীকরণ: $r = 6\cos\theta + 8\sin\theta$
    উভয়পক্ষকে $r$ দ্বারা গুণ করে কার্তেসীয় সমীকরণে রূপান্তর করি:
    \[ r^2 = 6r\cos\theta + 8r\sin\theta \implies x^2 + y^2 - 6x - 8y = 0 \]
    যেহেতু সমীকরণে কোনো ধ্রুবক পদ নেই (অর্থাৎ $c=0$), বৃত্তটি মূলবিন্দু $O(0, 0)$ (মেরু) দিয়ে অতিক্রম করে। $\theta = \frac{\pi}{4}$ রেখাটি মূলবিন্দুগামী একটি সরলরেখা। জ্যা-এর একটি প্রান্তবিন্দু অবশ্যই মূলবিন্দু $(0, 0)$। 
    
    \vspace{0.5em}
    জ্যা-এর দৈর্ঘ্য হবে মূলবিন্দু হতে অন্য ছেদবিন্দুর দূরত্ব, যা বৃত্তের সমীকরণে সরাসরি $\theta = \frac{\pi}{4}$ বসালে পাওয়া যাবে:
    \[ r = 6\cos\left(\frac{\pi}{4}\right) + 8\sin\left(\frac{\pi}{4}\right) \]
    \[ r = 6\left(\frac{1}{\sqrt{2}}\right) + 8\left(\frac{1}{\sqrt{2}}\right) = \frac{14}{\sqrt{2}} = 7\sqrt{2} \]
    অতএব, জ্যা-এর দৈর্ঘ্য $7\sqrt{2}$ একক।
\end{tcolorbox}

% ------------------------------------------
% QUESTION SECTION
% ------------------------------------------
\begin{tcolorbox}[questionbox={প্রশ্ন (Question) 5}]
    \noindent
    $x^2 + y^2 = 25$ এবং $x^2 + y^2 - 8x - 8y + 7 = 0$ বৃত্তদ্বয়ের সাধারণ জ্যা-কে ব্যাস ধরে অঙ্কিত বৃত্তটি $y$-অক্ষ থেকে কত একক দীর্ঘ জ্যা কর্তন করে?
    
    \vspace{0.8em}
    \begin{english}
    \noindent\textit{\color{darkgray}
    What is the length of the intercept made on the $y$-axis by the circle which has the common chord of the circles $x^2 + y^2 = 25$ and $x^2 + y^2 - 8x - 8y + 7 = 0$ as a diameter?
    }
    \end{english}
\end{tcolorbox}

% ------------------------------------------
% OPTIONS SECTION
% ------------------------------------------
\begin{itemize}[label={}, leftmargin=1cm, itemsep=0.5em]
    \item[\textbf{A)}] $2\sqrt{13}$ একক
    \item[\textbf{B)}] $4\sqrt{3}$ একক
    \item[\textbf{C)}] $5$ একক
    \item[\textbf{D)}] $2\sqrt{15}$ একক
\end{itemize}

% ------------------------------------------
% ANSWER HIGHLIGHT
% ------------------------------------------
\vspace{0.5em}
\begin{tcolorbox}[answerbox]
    \textbf{\color{success} উত্তর:} A) $2\sqrt{13}$ একক
\end{tcolorbox}
\vspace{1em}

% ------------------------------------------
% SOLUTION SECTION
% ------------------------------------------
\begin{tcolorbox}[solutionbox={সমাধান (Solution)}]
    \noindent
    বৃত্তদ্বয়:
    \begin{align*}
        S_1 &= x^2 + y^2 - 25 = 0 \\
        S_2 &= x^2 + y^2 - 8x - 8y + 7 = 0
    \end{align*}
    বৃত্তদ্বয়ের সাধারণ জ্যা-এর সমীকরণ $L = S_1 - S_2 = 0$:
    \[ (x^2 + y^2 - 25) - (x^2 + y^2 - 8x - 8y + 7) = 0 \]
    \[ 8x + 8y - 32 = 0 \implies x + y - 4 = 0 \]
    ছেদবিন্দুগামী যেকোনো বৃত্তের সমীকরণ $S_1 + \lambda L = 0$:
    \[ x^2 + y^2 - 25 + \lambda(x + y - 4) = 0 \]
    \[ x^2 + y^2 + \lambda x + \lambda y - (4\lambda + 25) = 0 \]
    বৃত্তটির কেন্দ্র $C\left(-\frac{\lambda}{2}, -\frac{\lambda}{2}\right)$। যেহেতু জ্যা-টি এই বৃত্তের ব্যাস, তাই বৃত্তের কেন্দ্রটি জ্যা-এর সমীকরণ $x + y - 4 = 0$ কে সিদ্ধ করবে:
    \[ -\frac{\lambda}{2} - \frac{\lambda}{2} - 4 = 0 \implies -\lambda = 4 \implies \lambda = -4 \]
    $\lambda = -4$ বৃত্তের সমীকরণে বসিয়ে পাই:
    \[ x^2 + y^2 - 4x - 4y - (4(-4) + 25) = 0 \]
    \[ x^2 + y^2 - 4x - 4y - 9 = 0 \]
    এই বৃত্তের কোটি প্যারামিটার $f = -2$ এবং ধ্রুবক $c = -9$। 
    $y$-অক্ষ হতে কর্তিত জ্যা-এর দৈর্ঘ্য:
    \[ 2\sqrt{f^2 - c} = 2\sqrt{(-2)^2 - (-9)} = 2\sqrt{4 + 9} = 2\sqrt{13} \]
    অতএব, কর্তিত জ্যা-এর দৈর্ঘ্য $2\sqrt{13}$ একক।
\end{tcolorbox}

% ------------------------------------------
% QUESTION SECTION
% ------------------------------------------
\begin{tcolorbox}[questionbox={প্রশ্ন (Question) 6}]
    \noindent
    $\theta \in [0, 2\pi]$ এর কতটি মানের জন্য সরলরেখা $x \sin \theta - y \cos \theta = 1$ বৃত্ত $x^2 + y^2 - 2x - 2y + 1 = 0$ এর একটি স্পর্শক হবে?
    
    \vspace{0.8em}
    \begin{english}
    \noindent\textit{\color{darkgray}
    For how many values of $\theta \in [0, 2\pi]$ is the line $x \sin \theta - y \cos \theta = 1$ a tangent to the circle $x^2 + y^2 - 2x - 2y + 1 = 0$?
    }
    \end{english}
\end{tcolorbox}

% ------------------------------------------
% OPTIONS SECTION
% ------------------------------------------
\begin{itemize}[label={}, leftmargin=1cm, itemsep=0.5em]
    \item[\textbf{A)}] $1$
    \item[\textbf{B)}] $2$
    \item[\textbf{C)}] $3$
    \item[\textbf{D)}] $4$
\end{itemize}

% ------------------------------------------
% ANSWER HIGHLIGHT
% ------------------------------------------
\vspace{0.5em}
\begin{tcolorbox}[answerbox]
    \textbf{\color{success} উত্তর:} B) $2$
\end{tcolorbox}
\vspace{1em}

% ------------------------------------------
% SOLUTION SECTION
% ------------------------------------------
\begin{tcolorbox}[solutionbox={সমাধান (Solution)}]
    \noindent
    \textbf{ধাপ ১: বৃত্তের কেন্দ্র ও ব্যাসার্ধ নির্ণয়} \\
    প্রদত্ত বৃত্তের সমীকরণ: $x^2 + y^2 - 2x - 2y + 1 = 0$
    \[ \implies (x - 1)^2 + (y - 1)^2 = 1^2 \]
    সুতরাং, বৃত্তের কেন্দ্র $C(1, 1)$ এবং ব্যাসার্ধ $R = 1$।
    
    \vspace{0.5em}
    \noindent\textbf{ধাপ ২: স্পর্শকের শর্ত প্রয়োগ} \\
    স্পর্শক হওয়ার শর্তানুযায়ী কেন্দ্র হতে রেখার লম্ব দূরত্ব ব্যাসার্ধের সমান হবে:
    \[ \frac{|\sin\theta - \cos\theta - 1|}{\sqrt{\sin^2\theta + \cos^2\theta}} = 1 \implies |\sin\theta - \cos\theta - 1| = 1 \]
    
    \vspace{0.5em}
    \noindent\textbf{ধাপ ৩: সমীকরণ সমাধান} \\
    মডুলাস উঠিয়ে দুটি কেস বিবেচনা করি:
    \begin{itemize}
        \item \textbf{প্রথম কেস:} 
        \[ \sin\theta - \cos\theta - 1 = 1 \implies \sin\theta - \cos\theta = 2 \]
        যা অসম্ভব, কারণ $\sin\theta - \cos\theta$ এর সর্বোচ্চ মান $\sqrt{2}$।
        \item \textbf{দ্বিতীয় কেস:} 
        \[ \sin\theta - \cos\theta - 1 = -1 \implies \sin\theta - \cos\theta = 0 \implies \sin\theta = \cos\theta \implies \tan\theta = 1 \]
    \end{itemize}
    
    \vspace{0.5em}
    \noindent\textbf{ধাপ ৪: মানের সীমা নির্ধারণ} \\
    $[0, 2\pi]$ সীমার মধ্যে $\tan\theta = 1$ এর দুটি সমাধান হলো:
    \[ \theta = \frac{\pi}{4} \quad \text{এবং} \quad \theta = \frac{5\pi}{4} \]
    অতএব, $\theta$ এর মান ২টি। সঠিক উত্তর B।
\end{tcolorbox}

% ------------------------------------------
% QUESTION SECTION
% ------------------------------------------
\begin{tcolorbox}[questionbox={প্রশ্ন (Question) 7}]
    \noindent
    একটি বৃত্ত $S$ মূলবিন্দু $(0,0)$ এবং $(1, 1)$ বিন্দুদ্বয় দিয়ে অতিক্রম করে। বৃত্তটি $x^2 + y^2 - 4x + 3 = 0$ বৃত্তকে লম্বভাবে ছেদ (orthogonally cuts) করে। বৃত্ত $S$ দ্বারা $y$-অক্ষের খণ্ডিত অংশের দৈর্ঘ্য কত একক?
    
    \vspace{0.8em}
    \begin{english}
    \noindent\textit{\color{darkgray}
    Let $S$ be a circle passing through $(0,0)$ and $(1, 1)$ and cutting the circle $x^2 + y^2 - 4x + 3 = 0$ orthogonally. What is the length of the y-intercept of the circle $S$?
    }
    \end{english}
\end{tcolorbox}

% ------------------------------------------
% OPTIONS SECTION
% ------------------------------------------
\begin{itemize}[label={}, leftmargin=1cm, itemsep=0.5em]
    \item[\textbf{A)}] $0.5$ একক
    \item[\textbf{B)}] $1.5$ একক
    \item[\textbf{C)}] $1.0$ একক
    \item[\textbf{D)}] $2.0$ একক
\end{itemize}

% ------------------------------------------
% ANSWER HIGHLIGHT
% ------------------------------------------
\vspace{0.5em}
\begin{tcolorbox}[answerbox]
    \textbf{\color{success} উত্তর:} A) $0.5$ একক
\end{tcolorbox}
\vspace{1em}

% ------------------------------------------
% SOLUTION SECTION
% ------------------------------------------
\begin{tcolorbox}[solutionbox={সমাধান (Solution)}]
    \noindent
    \textbf{ধাপ ১: বৃত্তের আদর্শ সমীকরণ গঠন} \\
    ধরি, বৃত্তের সমীকরণ: 
    \[ x^2 + y^2 + 2gx + 2fy + c = 0 \]
    
    \vspace{0.5em}
    \noindent\textbf{ধাপ ২: মূলবিন্দুগামী হওয়ার শর্ত প্রয়োগ} \\
    যেহেতু বৃত্তটি মূলবিন্দু $(0,0)$ দিয়ে অতিক্রম করে, তাই সমীকরণে $(0,0)$ বসিয়ে পাই:
    \[ 0^2 + 0^2 + 2g(0) + 2f(0) + c = 0 \implies c = 0 \]
    
    \vspace{0.5em}
    \noindent\textbf{ধাপ ৩: $(1, 1)$ বিন্দুগামী হওয়ার শর্ত প্রয়োগ} \\
    বৃত্তটি $(1, 1)$ বিন্দু দিয়ে অতিক্রম করায়:
    \[ 1^2 + 1^2 + 2g(1) + 2f(1) + 0 = 0 \implies 2g + 2f + 2 = 0 \implies g + f = -1 \quad \text{--- (i)} \]
    
    \vspace{0.5em}
    \noindent\textbf{ধাপ ৪: লম্বভাবে ছেদ করার শর্ত প্রয়োগ} \\
    প্রদত্ত দ্বিতীয় বৃত্ত হলো $x^2 + y^2 - 4x + 3 = 0$। এর ক্ষেত্রে $g_2 = -2$, $f_2 = 0$ এবং $c_2 = 3$।
    বৃত্ত দুটি লম্বভাবে ছেদ করার শর্ত ($2g_1g_2 + 2f_1f_2 = c_1 + c_2$) অনুযায়ী:
    \[ 2g(-2) + 2f(0) = 0 + 3 \implies -4g = 3 \implies g = -\frac{3}{4} \]
    
    \vspace{0.5em}
    \noindent\textbf{ধাপ ৫: $f$ এর মান নির্ণয় এবং বৃত্তের সমীকরণ গঠন} \\
    সমীকরণ (i) হতে $f$ এর মান পাই:
    \[ f = -1 - g = -1 - \left(-\frac{3}{4}\right) = -1 + \frac{3}{4} = -\frac{1}{4} \]
    অতএব, নির্ণেয় বৃত্ত $S$ এর সমীকরণ:
    \[ x^2 + y^2 -\frac{3}{2}x - \frac{1}{2}y = 0 \]
    
    \vspace{0.5em}
    \noindent\textbf{ধাপ ৬: $y$-অক্ষের খণ্ডিত অংশের দৈর্ঘ্য নির্ণয়} \\
    বৃত্তটি দ্বারা $y$-অক্ষের খণ্ডিত অংশ পেতে $x = 0$ বসাই:
    \[ 0 + y^2 - 0 - \frac{1}{2}y = 0 \implies y^2 - \frac{1}{2}y = 0 \implies y\left(y - \frac{1}{2}\right) = 0 \]
    এখান থেকে পাই $y = 0$ এবং পোষক বিন্দু $y = 0.5$। 
    অতএব, $y$-অক্ষের খণ্ডিত অংশের দৈর্ঘ্য $|0.5 - 0| = 0.5$ একক।
    অতএব, সঠিক উত্তর A।
\end{tcolorbox}

% ------------------------------------------
% QUESTION SECTION
% ------------------------------------------
\begin{tcolorbox}[questionbox={প্রশ্ন (Question) 8}]
    \noindent
    $x^2 + y^2 = 16$ বৃত্তের একটি জ্যা-এর সমীকরণ $x\cos\alpha + y\sin\alpha = p$। যদি জ্যা-টি বৃত্তের কেন্দ্রে $90^{\circ}$ কোণ উৎপন্ন করে, তবে $p$ এর মান কত একক?
    
    \vspace{0.8em}
    \begin{english}
    \noindent\textit{\color{darkgray}
    The equation of a chord of the circle $x^2 + y^2 = 16$ is $x\cos\alpha + y\sin\alpha = p$. If this chord subtends an angle of $90^{\circ}$ at the center of the circle, what is the value of $p$ in units?
    }
    \end{english}
\end{tcolorbox}

% ------------------------------------------
% OPTIONS SECTION
% ------------------------------------------
\begin{itemize}[label={}, leftmargin=1cm, itemsep=0.5em]
    \item[\textbf{A)}] $2\sqrt{2}$ একক
    \item[\textbf{B)}] $4\sqrt{2}$ একক
    \item[\textbf{C)}] $2$ একক
    \item[\textbf{D)}] $\sqrt{2}$ একক
\end{itemize}

% ------------------------------------------
% ANSWER HIGHLIGHT
% ------------------------------------------
\vspace{0.5em}
\begin{tcolorbox}[answerbox]
    \textbf{\color{success} উত্তর:} A) $2\sqrt{2}$ একক
\end{tcolorbox}
\vspace{1em}

% ------------------------------------------
% SOLUTION SECTION
% ------------------------------------------
\begin{tcolorbox}[solutionbox={সমাধান (Solution)}]
    \noindent
    বৃত্তের ব্যাসার্ধ $R = \sqrt{16} = 4$ একক এবং কেন্দ্র মূলবিন্দু $O(0, 0)$। কেন্দ্র $O(0, 0)$ হতে জ্যা $x\cos\alpha + y\sin\alpha - p = 0$ এর লম্ব দূরত্ব $d$:
    \[ d = \frac{|0\cos\alpha + 0\sin\alpha - p|}{\sqrt{\cos^2\alpha + \sin^2\alpha}} = |p| = p \quad (\text{ধরি } p > 0) \]
    যেহেতু জ্যা-টি কেন্দ্রে $90^{\circ}$ কোণ উৎপন্ন করে, তাই কেন্দ্র এবং জ্যা-এর প্রান্তবিন্দুদ্বয় মিলে একটি সমকোণী সমদ্বিবাহু ত্রিভুজ গঠন করে। ত্রিভুজটির জ্যামিতিক বৈশিষ্ট্য হতে আমরা পাই:
    \[ d = R \cos 45^{\circ} \]
    \[ p = 4 \times \frac{1}{\sqrt{2}} = 2\sqrt{2} \]
    অতএব, $p$ এর মান $2\sqrt{2}$ একক।
\end{tcolorbox}

% ------------------------------------------
% QUESTION SECTION
% ------------------------------------------
\begin{tcolorbox}[questionbox={প্রশ্ন (Question) 9}]
    \noindent
    বৃত্তদ্বয় $x^2 + y^2 - 4x - 6y - 12 = 0$ এবং $x^2 + y^2 + 2x + 2y - 2 = 0$ এর সাধারণ জ্যা-এর দৈর্ঘ্য কত একক?
    
    \vspace{0.8em}
    \begin{english}
    \noindent\textit{\color{darkgray}
    What is the length of the common chord of the circles $x^2 + y^2 - 4x - 6y - 12 = 0$ and $x^2 + y^2 + 2x + 2y - 2 = 0$ in units?
    }
    \end{english}
\end{tcolorbox}

% ------------------------------------------
% OPTIONS SECTION
% ------------------------------------------
\begin{itemize}[label={}, leftmargin=1cm, itemsep=0.5em]
    \item[\textbf{A)}] $\frac{8\sqrt{6}}{5}$ একক
    \item[\textbf{B)}] $\frac{4\sqrt{6}}{5}$ একক
    \item[\textbf{C)}] $\frac{6\sqrt{8}}{5}$ একক
    \item[\textbf{D)}] $4\sqrt{2}$ একক
\end{itemize}

% ------------------------------------------
% ANSWER HIGHLIGHT
% ------------------------------------------
\vspace{0.5em}
\begin{tcolorbox}[answerbox]
    \textbf{\color{success} উত্তর:} A) $\frac{8\sqrt{6}}{5}$ একক
\end{tcolorbox}
\vspace{1em}

% ------------------------------------------
% SOLUTION SECTION
% ------------------------------------------
\begin{tcolorbox}[solutionbox={সমাধান (Solution)}]
    \noindent
    বৃত্তদ্বয়ের সাধারণ জ্যা এর সমীকরণ $S_1 - S_2 = 0$:
    \[ (x^2 + y^2 - 4x - 6y - 12) - (x^2 + y^2 + 2x + 2y - 2) = 0 \]
    \[ -6x - 8y - 10 = 0 \implies 3x + 4y + 5 = 0 \]
    প্রথম বৃত্তের কেন্দ্র $C_1(2, 3)$ এবং ব্যাসার্ধ $r_1 = \sqrt{(-2)^2 + (-3)^2 - (-12)} = \sqrt{25} = 5$। 
    
    \vspace{0.5em}
    কেন্দ্র $C_1(2, 3)$ হতে জ্যা $3x + 4y + 5 = 0$ এর লম্ব দূরত্ব $d_1$:
    \[ d_1 = \frac{|3(2) + 4(3) + 5|}{\sqrt{3^2 + 4^2}} = \frac{|6 + 12 + 5}{5} = \frac{23}{5} = 4.6 \]
    সাধারণ জ্যা-এর সম্পূর্ণ দৈর্ঘ্য $L$:
    \[ L = 2\sqrt{r_1^2 - d_1^2} = 2\sqrt{5^2 - \left(\frac{23}{5}\right)^2} = 2\sqrt{25 - \frac{529}{25}} \]
    \[ = 2\sqrt{\frac{625 - 529}{25}} = 2\sqrt{\frac{96}{25}} = \frac{2 \times 4\sqrt{6}}{5} = \frac{8\sqrt{6}}{5} \]
    অতএব, সাধারণ জ্যা-এর দৈর্ঘ্য $\frac{8\sqrt{6}}{5}$ একক।
\end{tcolorbox}

% ------------------------------------------
% QUESTION SECTION
% ------------------------------------------
\begin{tcolorbox}[questionbox={প্রশ্ন (Question) 10}]
    \noindent
    $P(3, 3)$ বিন্দু হতে $x^2 + y^2 - 2x - 2y + 1 = 0$ বৃত্তে অঙ্কিত স্পর্শকদ্বয় এবং তাদের স্পর্শ জ্যা (chord of contact) দ্বারা গঠিত ত্রিভুজের ক্ষেত্রফল কত বর্গ একক?
    
    \vspace{0.8em}
    \begin{english}
    \noindent\textit{\color{darkgray}
    What is the area in square units of the triangle formed by the tangents drawn from the point $P(3, 3)$ to the circle $x^2 + y^2 - 2x - 2y + 1 = 0$ and their chord of contact?
    }
    \end{english}
\end{tcolorbox}

% ------------------------------------------
% OPTIONS SECTION
% ------------------------------------------
\begin{itemize}[label={}, leftmargin=1cm, itemsep=0.5em]
    \item[\textbf{A)}] $\frac{7\sqrt{7}}{8}$ বর্গ একক
    \item[\textbf{B)}] $\frac{3\sqrt{7}}{8}$ বর্গ একক
    \item[\textbf{C)}] $\frac{7\sqrt{3}}{4}$ বর্গ একক
    \item[\textbf{D)}] $\frac{\sqrt{7}}{2}$ বর্গ একক
\end{itemize}

% ------------------------------------------
% ANSWER HIGHLIGHT
% ------------------------------------------
\vspace{0.5em}
\begin{tcolorbox}[answerbox]
    \textbf{\color{success} উত্তর:} A) $\frac{7\sqrt{7}}{8}$ বর্গ একক
\end{tcolorbox}
\vspace{1em}

% ------------------------------------------
% SOLUTION SECTION
% ------------------------------------------
\begin{tcolorbox}[solutionbox={সমাধান (Solution)}]
    \noindent
    প্রদত্ত বৃত্তের সমীকরণ: $x^2 + y^2 - 2x - 2y + 1 = 0$
    বৃত্তের কেন্দ্র $C(1, 1)$ এবং ব্যাসার্ধ $r = \sqrt{(-1)^2 + (-1)^2 - 1} = 1$।
    
    \vspace{0.5em}
    বহিঃস্থ বিন্দু $P(3, 3)$ হতে অঙ্কিত স্পর্শকের দৈর্ঘ্য $L = \sqrt{S_1}$:
    \[ L = \sqrt{3^2 + 3^2 - 2(3) - 2(3) + 1} = \sqrt{9 + 9 - 6 - 6 + 1} = \sqrt{7} \]
    স্পর্শকদ্বয় এবং স্পর্শ জ্যা দ্বারা গঠিত ত্রিভুজের ক্ষেত্রফল নির্ণয়ের সূত্র:
    \[ \text{Area} = \frac{r L^3}{r^2 + L^2} \]
    \[ \text{Area} = \frac{1 \times (\sqrt{7})^3}{1^2 + (\sqrt{7})^2} = \frac{7\sqrt{7}}{1 + 7} = \frac{7\sqrt{7}}{8} \]
    অতএব, গঠিত ত্রিভুজের ক্ষেত্রফল $\frac{7\sqrt{7}}{8}$ বর্গ একক।
\end{tcolorbox}

% ------------------------------------------
% QUESTION SECTION
% ------------------------------------------
\begin{tcolorbox}[questionbox={প্রশ্ন (Question) 11}]
    \noindent
    একটি সমাক্ষীয় বৃত্তশ্রেণীর সীমাবদ্ধ বিন্দুদ্বয় (Limiting points) যথাক্রমে $(1, 1)$ and $(2, 2)$। এই সমাক্ষ শ্রেণীর মূল অক্ষের (Radical axis) সমীকরণ কোনটি?
    
    \vspace{0.8em}
    \begin{english}
    \noindent\textit{\color{darkgray}
    The limiting points of a coaxial system of circles are $(1, 1)$ and $(2, 2)$ respectively. Which of the following is the equation of the radical axis of this system?
    }
    \end{english}
\end{tcolorbox}

% ------------------------------------------
% OPTIONS SECTION
% ------------------------------------------
\begin{itemize}[label={}, leftmargin=1cm, itemsep=0.5em]
    \item[\textbf{A)}] $x + y - 3 = 0$
    \item[\textbf{B)}] $x - y = 0$
    \item[\textbf{C)}] $x + y + 3 = 0$
    \item[\textbf{D)}] $2x + 2y - 5 = 0$
\end{itemize}

% ------------------------------------------
% ANSWER HIGHLIGHT
% ------------------------------------------
\vspace{0.5em}
\begin{tcolorbox}[answerbox]
    \textbf{\color{success} উত্তর:} A) $x + y - 3 = 0$
\end{tcolorbox}
\vspace{1em}

% ------------------------------------------
% SOLUTION SECTION
% ------------------------------------------
\begin{tcolorbox}[solutionbox={সমাধান (Solution)}]
    \noindent
    সমাক্ষীয় বৃত্তশ্রেণীর সীমাবদ্ধ বিন্দুদ্বয় মূলত শূন্য ব্যাসার্ধের দুটি বিন্দু বৃত্ত নির্দেশ করে। আমরা জানি, সমাক্ষীয় বৃত্তশ্রেণীর মূল অক্ষটি (Radical Axis) সীমাবদ্ধ বিন্দুদ্বয়ের সংযোজক সরলরেখাংশের লম্ব দ্বিখণ্ডক (perpendicular bisector) হয়। 
    
    \vspace{0.5em}
    সীাবদ্ধ বিন্দুদ্বয় $P(1, 1)$ এবং $Q(2, 2)$ এর মধ্যবিন্দু $M$:
    \[ M = \left(\frac{1+2}{2}, \frac{1+2}{2}\right) = (1.5, 1.5) \]
    $PQ$ রেখার ঢাল $m_1 = \frac{2 - 1}{2 - 1} = 1$। অতএব, লম্ব দ্বিখণ্ডকের (মূল অক্ষের) ঢাল $m_2 = -\frac{1}{m_1} = -1$।
    মূল অক্ষের সমীকরণ:
    \[ y - 1.5 = -1(x - 1.5) \]
    \[ y - 1.5 = -x + 1.5 \]
    \[ x + y - 3 = 0 \]
    অতএব, মূল অক্ষের সমীকরণ $x + y - 3 = 0$।
\end{tcolorbox}

% ------------------------------------------
% QUESTION SECTION
% ------------------------------------------
\begin{tcolorbox}[questionbox={প্রশ্ন (Question) 12}]
    \noindent
    ধরি, দুটি বৃত্ত $C_1$ ও $C_2$ পরস্পরকে $P(1, 1)$ বিন্দুতে বাহ্যিক স্পর্শ করে। বৃত্তদ্বয় একই সাথে $S: x^2 + y^2 - 4x - 4y + 7 = 0$ বৃত্তটিকে লম্বভাবে ছেদ (orthogonally intersect) করে। যদি $P$ বিন্দুতে বৃত্তদ্বয়ের সাধারণ স্পর্শকের (common tangent) সমীকরণ $y = mx + c$ হয়, তবে $m + c$ এর মান কত?
    
    \vspace{0.8em}
    \begin{english}
    \noindent\textit{\color{darkgray}
    Let two circles $C_1$ and $C_2$ touch each other externally at the point $P(1, 1)$. Both $C_1$ and $C_2$ intersect the circle $S: x^2 + y^2 - 4x - 4y + 7 = 0$ orthogonally. If the equation of the common tangent to $C_1$ and $C_2$ at $P$ is $y = mx + c$, then the value of $m + c$ is:
    }
    \end{english}
\end{tcolorbox}

% ------------------------------------------
% OPTIONS SECTION
% ------------------------------------------
\begin{itemize}[label={}, leftmargin=1cm, itemsep=0.5em]
    \item[\textbf{A)}] $1$
    \item[\textbf{B)}] $2$
    \item[\textbf{C)}] $0$
    \item[\textbf{D)}] $-1$
\end{itemize}

% ------------------------------------------
% ANSWER HIGHLIGHT
% ------------------------------------------
\vspace{0.5em}
\begin{tcolorbox}[answerbox]
    \textbf{\color{success} উত্তর:} A) $1$
\end{tcolorbox}
\vspace{1em}

% ------------------------------------------
% SOLUTION SECTION
% ------------------------------------------
\begin{tcolorbox}[solutionbox={সমাধান (Solution)}]
    \noindent
    \textbf{ধাপ ১: বৃত্তের কেন্দ্রের অবস্থান বিশ্লেষণ} \\
    প্রদত্ত বৃত্ত $S: x^2 + y^2 - 4x - 4y + 7 = 0$। 
    এর কেন্দ্র $M(2, 2)$ এবং ব্যাসার্ধ $R_0 = \sqrt{2^2 + 2^2 - 7} = \sqrt{4 + 4 - 7} = \sqrt{1} = 1$।
    
    \vspace{0.5em}
    \noindent\textbf{ধাপ ২: স্পর্শ বৃত্ত ও মূল অক্ষের ধর্ম} \\
    যেহেতু $C_1$ এবং $C_2$ বৃত্ত দুটি পরস্পরকে $P(1, 1)$ বিন্দুতে স্পর্শ করে, তাই তাদের স্পর্শবিন্দুতে অঙ্কিত সাধারণ স্পর্শকটিই হবে উক্ত বৃত্তদ্বয়ের মূল অক্ষ (Radical Axis)।
    
    \vspace{0.5em}
    \noindent\textbf{ধাপ ৩: লম্ব ছেদের জ্যামিতিক বৈশিষ্ট্য} \\
    যেহেতু $C_1$ এবং $C_2$ উভয় বৃত্তই $S$ বৃত্তকে লম্বভাবে ছেদ করে, তাই $S$ বৃত্তের কেন্দ্র $M(2, 2)$ হতে $C_1$ ও $C_2$ বৃত্তের সাপেক্ষে স্পর্শকের দৈর্ঘ্যের বর্গ সমান হবে। এর অর্থ হলো, $S$ বৃত্তের কেন্দ্র $M(2, 2)$ অবশ্যই $C_1$ ও $C_2$ এর মূল অক্ষের (অর্থাৎ সাধারণ স্পর্শকের) ওপর অবস্থিত হবে।
    
    \vspace{0.5em}
    \noindent\textbf{ধাপ ৪: সাধারণ স্পর্শকের সমীকরণ নির্ণয়} \\
    যেহেতু সাধারণ স্পর্শকটি স্পর্শবিন্দু $P(1, 1)$ এবং $S$ বৃত্তের কেন্দ্র $M(2, 2)$ উভয় বিন্দুগামী, তাই এর সমীকরণ হবে:
    \[ y - 1 = \frac{2 - 1}{2 - 1}(x - 1) \implies y - 1 = 1(x - 1) \implies y - 1 = x - 1 \implies y = x \]
    
    \vspace{0.5em}
    \noindent\textbf{ধাপ ৫: সহগ তুলনা ও মান নির্ণয়} \\
    সমীকরণটিকে $y = mx + c$ এর সাথে তুলনা করে পাই:
    \[ m = 1 \quad \text{এবং} \quad c = 0 \]
    অতএব, নির্ণেয় মান:
    \[ m + c = 1 + 0 = 1 \]
    অতএব, সঠিক উত্তর A।
\end{tcolorbox}
% ------------------------------------------
% QUESTION SECTION
% ------------------------------------------
\begin{tcolorbox}[questionbox={প্রশ্ন (Question) 13}]
    \noindent
    $x^2 + y^2 - 2x - 2y + 1 = 0$ বৃত্তের যেসব জ্যা মূলবিন্দু $(0, 0)$ দিয়ে যায়, তাদের মধ্যবিন্দুসমূহের সঞ্চারপথের সমীকরণ কোনটি?
    
    \vspace{0.8em}
    \begin{english}
    \noindent\textit{\color{darkgray}
    What is the equation of the locus of the midpoints of the chords of the circle $x^2 + y^2 - 2x - 2y + 1 = 0$ that pass through the origin $(0, 0)$?
    }
    \end{english}
\end{tcolorbox}

% ------------------------------------------
% OPTIONS SECTION
% ------------------------------------------
\begin{itemize}[label={}, leftmargin=1cm, itemsep=0.5em]
    \item[\textbf{A)}] $x^2 + y^2 - x - y = 0$
    \item[\textbf{B)}] $x^2 + y^2 - 2x - 2y = 0$
    \item[\textbf{C)}] $x^2 + y^2 - \frac{1}{2}x - \frac{1}{2}y = 0$
    \item[\textbf{D)}] $x^2 + y^2 - 2x - y = 0$
\end{itemize}

% ------------------------------------------
% ANSWER HIGHLIGHT
% ------------------------------------------
\vspace{0.5em}
\begin{tcolorbox}[answerbox]
    \textbf{\color{success} উত্তর:} A) $x^2 + y^2 - x - y = 0$
\end{tcolorbox}
\vspace{1em}

% ------------------------------------------
% SOLUTION SECTION
% ------------------------------------------
\begin{tcolorbox}[solutionbox={সমাধান (Solution)}]
    \noindent
    বৃত্তের সমীকরণ: $x^2 + y^2 - 2x - 2y + 1 = 0$ এর কেন্দ্র $C(1, 1)$। 
    
    \vspace{0.5em}
    আমরা জানি, বৃত্তের কোনো জ্যা যদি একটি নির্দিষ্ট বিন্দু $P(x_1, y_1)$ দিয়ে যায়, তবে তার মধ্যবিন্দুর সঞ্চারপথ একটি বৃত্ত নির্দেশ করে। এই সঞ্চারপথ বৃত্তটির ব্যাস হয় বৃত্তের কেন্দ্র $C$ এবং নির্দিষ্ট বিন্দু $P$ এর সংযোজক রেখাংশ। 
    এখানে, কেন্দ্র $C(1, 1)$ এবং নির্দিষ্ট বিন্দু হলো মূলবিন্দু $O(0, 0)$। $C(1, 1)$ এবং $O(0, 0)$ কে ব্যাসের প্রান্তবিন্দু ধরে অঙ্কিত বৃত্তের সমীকরণ:
    \[ (x - 1)(x - 0) + (y - 1)(y - 0) = 0 \]
    \[ x(x - 1) + y(y - 1) = 0 \]
    \[ x^2 + y^2 - x - y = 0 \]
    অতএব, জ্যা-সমূহের মধ্যবিন্দুসমূহের সঞ্চারপথের সমীকরণ $x^2 + y^2 - x - y = 0$।
\end{tcolorbox}

% ------------------------------------------
% QUESTION SECTION
% ------------------------------------------
\begin{tcolorbox}[questionbox={প্রশ্ন (Question) 14}]
    \noindent
    যদি বৃত্ত $x^2 + y^2 = R^2$ এবং উপবৃত্ত $9x^2 + 16y^2 = 144$ এর সাধারণ জ্যা-সমূহ উপবৃত্তের কেন্দ্রে (origin) সমকোণ উৎপন্ন করে, তবে $R^2$ এর মান কত?
    
    \vspace{0.8em}
    \begin{english}
    \noindent\textit{\color{darkgray}
    If the common chords of the circle $x^2 + y^2 = R^2$ and the ellipse $9x^2 + 16y^2 = 144$ subtend a right angle at the center of the ellipse, then $R^2$ is equal to:
    }
    \end{english}
\end{tcolorbox}

% ------------------------------------------
% OPTIONS SECTION
% ------------------------------------------
\begin{itemize}[label={}, leftmargin=1cm, itemsep=0.5em]
    \item[\textbf{A)}] $\frac{288}{25}$
    \item[\textbf{B)}] $\frac{144}{25}$
    \item[\textbf{C)}] $\frac{144}{7}$
    \item[\textbf{D)}] $\frac{288}{7}$
\end{itemize}

% ------------------------------------------
% ANSWER HIGHLIGHT
% ------------------------------------------
\vspace{0.5em}
\begin{tcolorbox}[answerbox]
    \textbf{\color{success} উত্তর:} A) $\frac{288}{25}$
\end{tcolorbox}
\vspace{1em}

% ------------------------------------------
% SOLUTION SECTION
% ------------------------------------------
\begin{tcolorbox}[solutionbox={সমাধান (Solution)}]
    \noindent
    \textbf{ধাপ ১: উপবৃত্তের আদর্শ সমীকরণ রূপান্তর} \\
    প্রদত্ত উপবৃত্তের সমীকরণ: $9x^2 + 16y^2 = 144$
    \[ \implies \frac{x^2}{16} + \frac{y^2}{9} = 1 \quad \text{--- (i)} \]
    এখানে বৃহৎ অক্ষের অর্ধ-দৈর্ঘ্য $a^2 = 16$ এবং ক্ষুদ্র অক্ষের অর্ধ-দৈর্ঘ্য $b^2 = 9$।
    
    \vspace{0.5em}
    \noindent\textbf{ধাপ ২: সমজাতীয়করণ (Homogenization) পদ্ধতি প্রয়োগ} \\
    বৃত্তের সমীকরণ হতে পাই: $x^2 + y^2 = R^2 \implies \frac{x^2 + y^2}{R^2} = 1$। 
    উপবৃত্তের কেন্দ্র (মূলবিন্দু) থেকে ছেদবিন্দুদ্বয়ের সংযোগকারী সরলরেখাদ্বয়ের সমীকরণ পেতে উপবৃত্তের সমীকরণকে বৃত্তের এই সমীকরণ দ্বারা সমজাতীয় করি:
    \[ \frac{x^2}{16} + \frac{y^2}{9} = 1 \times 1 \implies \frac{x^2}{16} + \frac{y^2}{9} = \frac{x^2 + y^2}{R^2} \]
    \[ \implies x^2 \left( \frac{1}{16} - \frac{1}{R^2} \right) + y^2 \left( \frac{1}{9} - \frac{1}{R^2} \right) = 0 \]
    
    \vspace{0.5em}
    \noindent\textbf{ধাপ ৩: সমকোণ উৎপন্ন করার শর্ত প্রয়োগ} \\
    যেহেতু মূলবিন্দুতে উৎপন্ন কোণটি সমকোণ (অর্থাৎ রেখাদ্বয় পরস্পর লম্ব), তাই দ্বিঘাত সমীকরণে $x^2$ এবং $y^2$ এর সহগদ্বয়ের সমষ্টি শূন্য হবে:
    \[ \left( \frac{1}{16} - \frac{1}{R^2} \right) + \left( \frac{1}{9} - \frac{1}{R^2} \right) = 0 \]
    \[ \implies \frac{1}{16} + \frac{1}{9} = \frac{2}{R^2} \]
    
    \vspace{0.5em}
    \noindent\textbf{ধাপ ৪: $R^2$ এর মান নির্ণয়} \\
    \[ \frac{9 + 16}{144} = \frac{2}{R^2} \implies \frac{25}{144} = \frac{2}{R^2} \]
    \[ \implies 25R^2 = 288 \implies R^2 = \frac{288}{25} \]
    অতএব, সঠিক উত্তর A।
\end{tcolorbox}
% ------------------------------------------
% QUESTION SECTION
% ------------------------------------------
\begin{tcolorbox}[questionbox={প্রশ্ন (Question) 15}]
    \noindent
    ধরি, $F$ হলো সমতলে অবস্থিত বৃত্তের একটি পরিবার যা $x^2 + y^2 - 2\lambda x + \lambda^2 - 1 = 0$ সমীকরণ দ্বারা সংজ্ঞায়িত, যেখানে $\lambda \in \mathbb{R}$। সমতলের যে অঞ্চল $S$ হতে এই পরিবারের অন্তত একটি বৃত্তে পরস্পর লম্ব স্পর্শক অঙ্কন করা সম্ভব, তা নিচের কোন অসমতা দ্বারা নির্দেশিত হয়?
    
    \vspace{0.8em}
    \begin{english}
    \noindent\textit{\color{darkgray}
    Let $F$ be the family of circles $x^2 + y^2 - 2\lambda x + \lambda^2 - 1 = 0$, where $\lambda \in \mathbb{R}$. Let $S$ be the region of all points $P(x, y)$ in the coordinate plane from which perpendicular tangents can be drawn to at least one circle of the family $F$. Then the region $S$ is bounded by:
    }
    \end{english}
\end{tcolorbox}

% ------------------------------------------
% OPTIONS SECTION
% ------------------------------------------
\begin{itemize}[label={}, leftmargin=1cm, itemsep=0.5em]
    \item[\textbf{A)}] $y^2 \le 2$
    \item[\textbf{B)}] $x^2 \le 2$
    \item[\textbf{C)}] $x^2 + y^2 \le 2$
    \item[\textbf{D)}] $x^2 - y^2 \le 2$
\end{itemize}

% ------------------------------------------
% ANSWER HIGHLIGHT
% ------------------------------------------
\vspace{0.5em}
\begin{tcolorbox}[answerbox]
    \textbf{\color{success} উত্তর:} A) $y^2 \le 2$
\end{tcolorbox}
\vspace{1em}

% ------------------------------------------
% SOLUTION SECTION
% ------------------------------------------
\begin{tcolorbox}[solutionbox={সমাধান (Solution)}]
    \noindent
    \textbf{ধাপ ১: বৃত্তের সাধারণ বৈশিষ্ট্য নির্ণয়} \\
    প্রদত্ত বৃত্তের সমীকরণ: 
    \[ x^2 + y^2 - 2\lambda x + \lambda^2 - 1 = 0 \implies (x - \lambda)^2 + y^2 = 1 \]
    এটি একটি বৃত্তের পরিবার যার কেন্দ্র $(\lambda, 0)$ এবং ব্যাসার্ধ $R = 1$ (যেখানে $\lambda$ একটি বাস্তব প্যারামিটার)।
    
    \vspace{0.5em}
    \noindent\textbf{ধাপ ২: নিয়ামক বৃত্ত (Director Circle) তত্ত্ব প্রয়োগ} \\
    যেকোনো বিন্দু থেকে বৃত্তে অঙ্কিত লম্ব স্পর্শকদ্বয়ের ছেদবিন্দুর সঞ্চারপথ হলো ওই বৃত্তের নিয়ামক বৃত্ত। নিয়ামক বৃত্তের সমীকরণ হবে কেন্দ্র $(\lambda, 0)$ এবং ব্যাসার্ধ $\sqrt{2}R = \sqrt{2}$ বিশিষ্ট বৃত্ত:
    \[ (x - \lambda)^2 + y^2 = (\sqrt{2})^2 \implies (x - \lambda)^2 + y^2 = 2 \]
    
    \vspace{0.5em}
    \noindent\textbf{ধাপ ৩: প্যারামিটার $\lambda$ এর বাস্তব মানের অস্তিত্ব} \\
    পরিবারের \textit{অন্তত একটি} বৃত্তের সাপেক্ষে নির্দিষ্ট বিন্দু $(x, y)$ হতে লম্ব স্পর্শক আঁকা সম্ভব হওয়ার অর্থ হলো সমীকরণটিতে বাস্তব $\lambda$ এর অস্তিত্ব থাকতে হবে। সমীকরণটিকে $\lambda$ এর সাপেক্ষে সাজাই:
    \[ (x - \lambda)^2 = 2 - y^2 \]
    যেহেতু বর্গের মান সর্বদা অ-ঋণাত্মক হয় (অর্থাৎ $(x - \lambda)^2 \ge 0$), তাই বাস্তব $\lambda$-এর অস্তিত্বের জন্য ডানপক্ষকেও অবশ্যই অ-ঋণাত্মক হতে হবে:
    \[ 2 - y^2 \ge 0 \implies y^2 \le 2 \]
    
    \vspace{0.5em}
    \noindent\textbf{ধাপ ৪: অঞ্চলের সীমানা নির্ধারণ} \\
    অতএব, নির্ণেয় অঞ্চলটি $y^2 \le 2$ অসমতা দ্বারা নির্দেশিত হয়।
    অতএব, সঠিক উত্তর A।
\end{tcolorbox}

% ------------------------------------------
% QUESTION SECTION
% ------------------------------------------
\begin{tcolorbox}[questionbox={প্রশ্ন (Question) 16}]
    \noindent
    $4x^2 + 4y^2 - 8x + 24y - 17 = 0$ বৃত্তের কেন্দ্র হতে $x - y - 6 = 0$ জ্যা-টির ওপর অঙ্কিত লম্বের পাদবিন্দু হতে জ্যা-টি $y$-অক্ষকে যে বিন্দুতে ছেদ করে তার দূরত্ব কত একক?
    
    \vspace{0.8em}
    \begin{english}
    \noindent\textit{\color{darkgray}
    What is the distance in units from the foot of the perpendicular drawn from the center of the circle $4x^2 + 4y^2 - 8x + 24y - 17 = 0$ to the chord $x - y - 6 = 0$ to the point where the chord intersects the $y$-axis?
    }
    \end{english}
\end{tcolorbox}

% ------------------------------------------
% OPTIONS SECTION
% ------------------------------------------
\begin{itemize}[label={}, leftmargin=1cm, itemsep=0.5em]
    \item[\textbf{A)}] $\sqrt{2}$ একক
    \item[\textbf{B)}] $2\sqrt{2}$ একক
    \item[\textbf{C)}] $4$ একক
    \item[\textbf{D)}] $3\sqrt{2}$ একক
\end{itemize}

% ------------------------------------------
% ANSWER HIGHLIGHT
% ------------------------------------------
\vspace{0.5em}
\begin{tcolorbox}[answerbox]
    \textbf{\color{success} উত্তর:} B) $2\sqrt{2}$ একক
\end{tcolorbox}
\vspace{1em}

% ------------------------------------------
% SOLUTION SECTION
% ------------------------------------------
\begin{tcolorbox}[solutionbox={সমাধান (Solution)}]
    \noindent
    1. প্রদত্ত বৃত্তের সমীকরণকে আদর্শ রূপ দেওয়ার জন্য $4$ দ্বারা ভাগ করে পাই:
    \[ x^2 + y^2 - 2x + 6y - \frac{17}{4} = 0 \]
    বৃত্তের কেন্দ্র $C(1, -3)$।
    
    \vspace{0.5em}
    2. জ্যা এর সমীকরণ: $x - y - 6 = 0$।
    
    3. কেন্দ্র $C(1, -3)$ হতে জ্যা-এর উপর লম্ব সরলরেখার সমীকরণ:
    \[ y - (-3) = -1(x - 1) \implies y + 3 = -x + 1 \implies x + y + 2 = 0 \]
    4. লম্বের পাদবিন্দুটি (foot of the perpendicular) হলো জ্যা এবং অঙ্কিত লম্ব রেখার ছেদবিন্দু:
    \[ (x - y - 6) + (x + y + 2) = 0 \implies 2x - 4 = 0 \implies x = 2 \]
    \[ y = x - 6 = 2 - 6 = -4 \]
    $\therefore$ পাদবিন্দুর স্থানাঙ্ক $P(2, -4)$।
    
    \vspace{0.5em}
    5. জ্যা-টি $y$-অক্ষকে স্পর্শ/ছেদ করে যে বিন্দুতে, তার ভুজ $x = 0$:
    \[ 0 - y - 6 = 0 \implies y = -6 \]
    $\therefore$ ছেদবিন্দুর স্থানাঙ্ক $Q(0, -6)$।
    
    6. পাদবিন্দু $P(2, -4)$ হতে ছেদবিন্দু $Q(0, -6)$ এর দূরত্ব:
    \[ PQ = \sqrt{(2 - 0)^2 + (-4 - (-6))^2} = \sqrt{2^2 + 2^2} = \sqrt{8} = 2\sqrt{2} \]
    অতএব, পাদবিন্দু হতে ছেদবিন্দুর দূরত্ব $2\sqrt{2}$ একক।
\end{tcolorbox}

% ------------------------------------------
% QUESTION SECTION
% ------------------------------------------
\begin{tcolorbox}[questionbox={প্রশ্ন (Question) 17}]
    \noindent
    $x^2 + y^2 - 6x - 8y + 5 = 0$ বৃত্তের অন্তর্লিখিত একটি সমবাহু ত্রিভুজের ক্ষেত্রফল কত বর্গ একক?
    
    \vspace{0.8em}
    \begin{english}
    \noindent\textit{\color{darkgray}
    What is the area in square units of an equilateral triangle inscribed in the circle $x^2 + y^2 - 6x - 8y + 5 = 0$?
    }
    \end{english}
\end{tcolorbox}

% ------------------------------------------
% OPTIONS SECTION
% ------------------------------------------
\begin{itemize}[label={}, leftmargin=1cm, itemsep=0.5em]
    \item[\textbf{A)}] $10\sqrt{3}$ বর্গ একক
    \item[\textbf{B)}] $12\sqrt{3}$ বর্গ একক
    \item[\textbf{C)}] $15\sqrt{3}$ বর্গ একক
    \item[\textbf{D)}] $20\sqrt{3}$ বর্গ একক
\end{itemize}

% ------------------------------------------
% ANSWER HIGHLIGHT
% ------------------------------------------
\vspace{0.5em}
\begin{tcolorbox}[answerbox]
    \textbf{\color{success} উত্তর:} C) $15\sqrt{3}$ বর্গ একক
\end{tcolorbox}
\vspace{1em}

% ------------------------------------------
% SOLUTION SECTION
% ------------------------------------------
\begin{tcolorbox}[solutionbox={সমাধান (Solution)}]
    \noindent
    1. প্রদত্ত বৃত্তের সমীকরণ: $x^2 + y^2 - 6x - 8y + 5 = 0$
    
    \vspace{0.5em}
    2. বৃত্তের ব্যাসার্ধ $R$:
    \[ R = \sqrt{(-3)^2 + (-4)^2 - 5} = \sqrt{9 + 16 - 5} = \sqrt{20} \]
    3. বৃত্তে অন্তর্লিখিত সমবাহু ত্রিভুজের বাহুর দৈর্ঘ্য $a$ এবং ব্যাসার্ধ $R$ এর সম্পর্ক:
    \[ a = R\sqrt{3} = \sqrt{20} \times \sqrt{3} = \sqrt{60} \]
    4. সমবাহু ত্রিভুজের ক্ষেত্রফল $A$ এর সূত্র:
    \[ A = \frac{\sqrt{3}}{4} a^2 = \frac{\sqrt{3}}{4} (\sqrt{60})^2 = \frac{\sqrt{3}}{4} \times 60 = 15\sqrt{3} \]
    অতএব, সমবাহু ত্রিভুজটির ক্ষেত্রফল $15\sqrt{3}$ বর্গ একক।
\end{tcolorbox}

% ------------------------------------------
% QUESTION SECTION
% ------------------------------------------
\begin{tcolorbox}[questionbox={প্রশ্ন (Question) 18}]
    \noindent
    একটি বৃত্ত মূলবিন্দু $O(0, 0)$ দিয়ে যায় এবং স্থানাঙ্কের ধনাত্মক অক্ষদ্বয় হতে যথাক্রমে $OA = 2a$ এবং $OB = 2b$ অংশ কর্তন করে। বৃত্তটির কেন্দ্র হতে জ্যা $AB$ এর লম্ব দূরত্ব কত একক?
    
    \vspace{0.8em}
    \begin{english}
    \noindent\textit{\color{darkgray}
    A circle passes through the origin $O(0, 0)$ and cuts off intercepts of lengths $OA = 2a$ and $OB = 2b$ from the positive coordinate axes respectively. What is the perpendicular distance in units from the center of the circle to the chord $AB$?
    }
    \end{english}
\end{tcolorbox}

% ------------------------------------------
% OPTIONS SECTION
% ------------------------------------------
\begin{itemize}[label={}, leftmargin=1cm, itemsep=0.5em]
    \item[\textbf{A)}] $a^2 + b^2$ একক
    \item[\textbf{B)}] $\frac{ab}{\sqrt{a^2 + b^2}}$ একক
    \item[\textbf{C)}] $0$ একক
    \item[\textbf{D)}] $\frac{a^2 + b^2}{2}$ একক
\end{itemize}

% ------------------------------------------
% ANSWER HIGHLIGHT
% ------------------------------------------
\vspace{0.5em}
\begin{tcolorbox}[answerbox]
    \textbf{\color{success} উত্তর:} C) $0$ একক
\end{tcolorbox}
\vspace{1em}

% ------------------------------------------
% SOLUTION SECTION
% ------------------------------------------
\begin{tcolorbox}[solutionbox={সমাধান (Solution)}]
    \noindent
    1. যেহেতু বৃত্তটি মূলবিন্দুগামী এবং অক্ষদ্বয় হতে যথাক্রমে $2a$ এবং $2b$ অংশ কর্তন করে, তাই অক্ষদ্বয়ের উপর ছেদবিন্দুদ্বয় যথাক্রমে $A(2a, 0)$ এবং $B(0, 2b)$।
    
    \vspace{0.5em}
    2. সমকোণী ত্রিভুজ $\triangle OAB$ এর পরিবৃত্তই হলো এই বৃত্তটি, কারণ $\angle AOB = 90^{\circ}$।
    
    3. বৃত্তের ব্যাস হলো অতিভুজ $AB$।
    
    4. বৃত্তের কেন্দ্র $C$ অবশ্যই ব্যাস $AB$ এর মধ্যবিন্দু হবে।
    
    5. যেহেতু বৃত্তের কেন্দ্রটি সরলরেখা $AB$ এর ওপরই অবস্থিত, সেহেতু কেন্দ্র হতে জ্যা $AB$ এর লম্ব দূরত্ব সর্বদা $0$ একক হবে।
\end{tcolorbox}

% ------------------------------------------
% QUESTION SECTION
% ------------------------------------------
\begin{tcolorbox}[questionbox={প্রশ্ন (Question) 19}]
    \noindent
    $(2, 3)$ এবং $(-1, 0)$ সীমাবদ্ধ বিন্দুদ্বয় (limiting points) বিশিষ্ট সমাক্ষীয় বৃত্তশ্রেণীর (coaxial system of circles) যে বৃত্তটি $2x + y = 3$ সরলরেখাকে স্পর্শ করে তার ব্যাসার্ধ কত একক?
    
    \vspace{0.8em}
    \begin{english}
    \noindent\textit{\color{darkgray}
    What is the radius in units of the circle of the coaxial system with limiting points $(2, 3)$ and $(-1, 0)$ that touches the line $2x + y = 3$?
    }
    \end{english}
\end{tcolorbox}

% ------------------------------------------
% OPTIONS SECTION
% ------------------------------------------
\begin{itemize}[label={}, leftmargin=1cm, itemsep=0.5em]
    \item[\textbf{A)}] $\frac{3}{\sqrt{5}}$ একক
    \item[\textbf{B)}] $\frac{3}{2}$ একক
    \item[\textbf{C)}] $2\sqrt{5}$ একক
    \item[\textbf{D)}] $3$ একক
\end{itemize}

% ------------------------------------------
% ANSWER HIGHLIGHT
% ------------------------------------------
\vspace{0.5em}
\begin{tcolorbox}[answerbox]
    \textbf{\color{success} উত্তর:} C) $2\sqrt{5}$ একক
\end{tcolorbox}
\vspace{1em}

% ------------------------------------------
% SOLUTION SECTION
% ------------------------------------------
\begin{tcolorbox}[solutionbox={সমাধান (Solution)}]
    \noindent
    1. সীমাবদ্ধ বিন্দুদ্বয় মূলত সমাক্ষীয় শ্রেণীর শূন্য ব্যাসার্ধের বৃত্তদ্বয় নির্দেশ করে:
    \begin{align*}
        S_1 &= (x - 2)^2 + (y - 3)^2 = x^2 + y^2 - 4x - 6y + 13 = 0 \\
        S_2 &= (x + 1)^2 + y^2 = x^2 + y^2 + 2x + 1 = 0
    \end{align*}
    2. সমাক্ষীয় শ্রেণীর মূল অক্ষ (Radical Axis) কাব্যিক নিয়মে $L = S_1 - S_2 = 0$:
    \[ -6x - 6y + 12 = 0 \implies x + y - 2 = 0 \]
    3. সমাক্ষীয় শ্রেণীর যেকোনো বৃত্তের সাধারণ সমীকরণ:
    \[ (x + 1)^2 + y^2 + \lambda(x + y - 2) = 0 \]
    \[ \implies x^2 + y^2 + (\lambda + 2)x + \lambda y + (1 - 2\lambda) = 0 \]
    বৃত্তটির কেন্দ্র $C\left(-\frac{\lambda + 2}{2}, -\frac{\lambda}{2}\right)$ এবং ব্যাসার্ধ $R$:
    \[ R = \sqrt{\left(\frac{\lambda + 2}{2}\right)^2 + \left(\frac{\lambda}{2}\right)^2 - (1 - 2\lambda)} = \sqrt{\frac{2\lambda^2 + 12\lambda}{4}} \]
    4. বৃত্তটি $2x + y - 3 = 0$ রেখাকে স্পর্শ করার শর্তানুযায়ী, কেন্দ্র হতে স্পর্শকের লম্ব দূরত্ব ব্যাসার্ধ $R$ এর সমান হবে:
    \[ \frac{\left|2\left(-\frac{\lambda + 2}{2}\right) + \left(-\frac{\lambda}{2}\right) - 3\right|}{\sqrt{2^2 + 1^2}} = R \implies \frac{|3\lambda + 10|}{2\sqrt{5}} = R \]
    5. বর্গ করে সমাধান করে পাই:
    \[ \frac{(3\lambda + 10)^2}{20} = \frac{2\lambda^2 + 12\lambda}{4} \implies 9\lambda^2 + 60\lambda + 100 = 10\lambda^2 + 60\lambda \]
    \[ \implies \lambda^2 = 100 \implies \lambda = -10 \quad (\text{ক্ষুদ্রতম স্পর্শক বৃত্তের জন্য}) \]
    6. $\lambda = -10$ হলে ব্যাসার্ধ $R$:
    \[ R = \sqrt{\frac{2(100) - 120}{4}} = \sqrt{20} = 2\sqrt{5} \]
    অতএব, স্পর্শকারী বৃত্তটির ব্যাসার্ধ $2\sqrt{5}$ একক।
\end{tcolorbox}

% ------------------------------------------
% QUESTION SECTION
% ------------------------------------------
\begin{tcolorbox}[questionbox={প্রশ্ন (Question) 20}]
    \noindent
    বৃত্তদ্বয় $C_1: x^2 + y^2 = 1$ এবং $C_2: x^2 + y^2 - 6x - 8y + 24 = 0$ বিবেচনা করো। একটি পরিবর্তনশীল বৃত্ত $C_3$ সর্বদা $C_1$ কে বাহ্যিক স্পর্শ করে এবং $C_2$ কে অন্তস্থভাবে স্পর্শ করে। $C_3$ এর কেন্দ্রের সঞ্চারপথ একটি অধিবৃত্ত (hyperbola) নির্দেশ করলে, উক্ত অধিবৃত্তের উৎকেন্দ্রিকতা (eccentricity) কত?
    
    \vspace{0.8em}
    \begin{english}
    \noindent\textit{\color{darkgray}
    Consider two circles $C_1: x^2 + y^2 = 1$ and $C_2: x^2 + y^2 - 6x - 8y + 24 = 0$. A variable circle $C_3$ touches $C_1$ externally and $C_2$ internally. If the locus of the center of $C_3$ is a hyperbola, then the eccentricity of this hyperbola is:
    }
    \end{english}
\end{tcolorbox}

% ------------------------------------------
% OPTIONS SECTION
% ------------------------------------------
\begin{itemize}[label={}, leftmargin=1cm, itemsep=0.5em]
    \item[\textbf{A)}] $2.5$
    \item[\textbf{B)}] $2.0$
    \item[\textbf{C)}] $1.5$
    \item[\textbf{D)}] $\sqrt{5}$
\end{itemize}

% ------------------------------------------
% ANSWER HIGHLIGHT
% ------------------------------------------
\vspace{0.5em}
\begin{tcolorbox}[answerbox]
    \textbf{\color{success} উত্তর:} A) $2.5$
\end{tcolorbox}
\vspace{1em}

% ------------------------------------------
% SOLUTION SECTION
% ------------------------------------------
\begin{tcolorbox}[solutionbox={সমাধান (Solution)}]
    \noindent
    \textbf{ধাপ ১: বৃত্তদ্বয়ের কেন্দ্র ও ব্যাসার্ধ নির্ণয়} \\
    * প্রথম বৃত্ত $C_1: x^2 + y^2 = 1$ এর কেন্দ্র $O(0, 0)$ এবং ব্যাসার্ধ $r_1 = 1$।
    * দ্বিতীয় বৃত্ত $C_2: x^2 + y^2 - 6x - 8y + 24 = 0$
    \[ \implies (x - 3)^2 + (y - 4)^2 = 3^2 + 4^2 - 24 = 9 + 16 - 24 = 1 \]
    সুতরাং, দ্বিতীয় বৃত্তের কেন্দ্র $A(3, 4)$ এবং ব্যাসার্ধ বৃত্তের জন্য $r_2 = \sqrt{1} = 1$।
    
    \vspace{0.5em}
    \noindent\textbf{ধাপ ২: স্পর্শের শর্ত প্রয়োগ} \\
    ধরি, পরিবর্তনশীল বৃত্ত $C_3$ এর কেন্দ্র $P(x, y)$ এবং ব্যাসার্ধ $r$।
    * $C_3$ বৃত্তটি $C_1$ কে বাহ্যিক স্পর্শ করার শর্তানুযায়ী:
    \[ PO = r + r_1 \implies PO = r + 1 \quad \text{--- (i)} \]
    * $C_3$ বৃত্তটি $C_2$ কে অন্তস্থভাবে স্পর্শ করার শর্তানুযায়ী:
    \[ PA = r_2 - r \implies PA = 1 - r \quad \text{--- (ii)} \]
    
    \vspace{0.5em}
    \noindent\textbf{ধাপ ৩: দূরত্বের অন্তরফল ও অধিবৃত্তের সমীকরণ গঠন} \\
    সমীকরণ (i) ও (ii) যোগ করে পাই:
    \[ PO + PA = (r + 1) + (1 - r) = 2 \]
    (অথবা বিয়োগ করলে দূরত্বের ধ্রুবক অন্তরফল পাওয়া যায়, যা অধিবৃত্তের শর্ত নির্দেশ করে)। এখানে কেন্দ্রদ্বয় থেকে দূরত্বের অন্তরফল:
    \[ |PO - PA| = |(r + 1) - (r - 1)| = 2 \]
    যেহেতু নির্দিষ্ট দুটি বিন্দু ($O$ এবং $A$) হতে দূরত্বের অন্তরফল সর্বদা একটি ধ্রুবক ($2$), তাই কেন্দ্র $P(x, y)$ এর সঞ্চারপথ একটি অধিবৃত্ত নির্দেশ করে যার নাভীদ্বয় $O(0,0)$ এবং $A(3,4)$।
    
    \vspace{0.5em}
    \noindent\textbf{ধাপ ৪: উৎকেন্দ্রিকতা ($e$) নির্ণয়} \\
    অধিবৃত্তের ক্ষেত্রে নাভীদ্বয়ের মধ্যবর্তী দূরত্ব $2ae$ এবং শীর্ষদ্বয়ের মধ্যবর্তী দূরত্ব $2a$:
    * নাভীদ্বয়ের মধ্যবর্তী দূরত্ব $2ae = OA$:
    \[ OA = \sqrt{(3 - 0)^2 + (4 - 0)^2} = \sqrt{3^2 + 4^2} = \sqrt{25} = 5 \]
    * দূরত্বের ধ্রুবক অন্তরফল $2a = 2 \implies a = 1$।
    
    অতএব, উৎকেন্দ্রিকতা $e$:
    \[ 2(1)e = 5 \implies e = \frac{5}{2} = 2.5 \]
    অতএব, সঠিক উত্তর A।
\end{tcolorbox}
% ------------------------------------------
% QUESTION SECTION
% ------------------------------------------
\begin{tcolorbox}[questionbox={প্রশ্ন (Question) 21}]
    \noindent
    জটিল তলে $|z - (1 + i)| = 2$ বৃত্তটির কেন্দ্র এবং ব্যাসার্ধ যথাক্রমে কত?
    
    \vspace{0.8em}
    \begin{english}
    \noindent\textit{\color{darkgray}
    What are the center and radius of the circle represented by $|z - (1 + i)| = 2$ in the complex plane?
    }
    \end{english}
\end{tcolorbox}

% ------------------------------------------
% OPTIONS SECTION
% ------------------------------------------
\begin{itemize}[label={}, leftmargin=1cm, itemsep=0.5em]
    \item[\textbf{A)}] কেন্দ্র $(1, -1)$, ব্যাসার্ধ $2$ একক
    \item[\textbf{B)}] কেন্দ্র $(1, 1)$, ব্যাসার্ধ $4$ একক
    \item[\textbf{C)}] কেন্দ্র $(1, 1)$, ব্যাসার্ধ $2$ একক
    \item[\textbf{D)}] কেন্দ্র $(-1, -1)$, ব্যাসার্ধ $2$ একক
\end{itemize}

% ------------------------------------------
% ANSWER HIGHLIGHT
% ------------------------------------------
\vspace{0.5em}
\begin{tcolorbox}[answerbox]
    \textbf{\color{success} উত্তর:} C) কেন্দ্র $(1, 1)$, ব্যাসার্ধ $2$ একক
\end{tcolorbox}
\vspace{1em}

% ------------------------------------------
% SOLUTION SECTION
% ------------------------------------------
\begin{tcolorbox}[solutionbox={সমাধান (Solution)}]
    \noindent
    1. ধরি, জটিল সংখ্যা $z = x + iy$।
    
    \vspace{0.5em}
    2. প্রদত্ত সম্পর্ক: $|z - (1 + i)| = 2$
    \[ \implies |(x + iy) - (1 + i)| = 2 \]
    \[ \implies |(x - 1) + i(y - 1)| = 2 \]
    3. জটিল সংখ্যার পরম মানের সংজ্ঞানুসারে:
    \[ \sqrt{(x - 1)^2 + (y - 1)^2} = 2 \]
    4. উভয়পক্ষকে বর্গ করে পাই:
    \[ (x - 1)^2 + (y - 1)^2 = 2^2 \]
    এটি কার্তেসীয় সমতলে একটি বৃত্ত নির্দেশ করে যার কেন্দ্র $(1, 1)$ এবং ব্যাসার্ধ $2$ একক।
\end{tcolorbox}

% ------------------------------------------
% QUESTION SECTION
% ------------------------------------------
\begin{tcolorbox}[questionbox={প্রশ্ন (Question) 22}]
    \noindent
    $S_1 \equiv x^2 + y^2 - 4 = 0$ এবং $S_2 \equiv x^2 + y^2 - 6x - 8y + 16 = 0$ বৃত্তদ্বয় দ্বারা গঠিত সমাক্ষ বৃত্ত জোটের (coaxial system of circles) যে বৃত্তটি $S_3 \equiv x^2 + y^2 - 2x - 2y = 0$ বৃত্তকে লম্বভাবে ছেদ (orthogonally cuts) করে, তার ব্যাসার্ধ কত একক?
    
    \vspace{0.8em}
    \begin{english}
    \noindent\textit{\color{darkgray}
    A circle $C$ belongs to the coaxial system generated by $S_1 \equiv x^2 + y^2 - 4 = 0$ and $S_2 \equiv x^2 + y^2 - 6x - 8y + 16 = 0$. If the circle $C$ cuts the circle $S_3 \equiv x^2 + y^2 - 2x - 2y = 0$ orthogonally, then the radius of $C$ is:
    }
    \end{english}
\end{tcolorbox}

% ------------------------------------------
% OPTIONS SECTION
% ------------------------------------------
\begin{itemize}[label={}, leftmargin=1cm, itemsep=0.5em]
    \item[\textbf{A)}] $\frac{4}{3}$ একক
    \item[\textbf{B)}] $\frac{5}{3}$ একক
    \item[\textbf{C)}] $2$ একক
    \item[\textbf{D)}] $\frac{2}{3}$ একক
\end{itemize}

% ------------------------------------------
% ANSWER HIGHLIGHT
% ------------------------------------------
\vspace{0.5em}
\begin{tcolorbox}[answerbox]
    \textbf{\color{success} উত্তর:} A) $\frac{4}{3}$ একক
\end{tcolorbox}
\vspace{1em}

% ------------------------------------------
% SOLUTION SECTION
% ------------------------------------------
\begin{tcolorbox}[solutionbox={সমাধান (Solution)}]
    \noindent
    \textbf{ধাপ ১: সমাক্ষ বৃত্তশ্রেণীর (Coaxial System) সমীকরণ গঠন} \\
    প্রদত্ত বৃত্তদ্বয় $S_1 = 0$ এবং $S_2 = 0$ দ্বারা গঠিত সমাক্ষ বৃত্তশ্রেণীর সাধারণ সমীকরণ লেখা যায়:
    \[ S_1 + \lambda(S_1 - S_2) = 0 \]
    এখানে মূল অক্ষ (Radical Axis) $S_1 - S_2 = 0$:
    \[ (x^2 + y^2 - 4) - (x^2 + y^2 - 6x - 8y + 16) = 0 \implies 6x + 8y - 20 = 0 \implies 3x + 4y - 10 = 0 \]
    অতএব, সমাক্ষ বৃত্তশ্রেণীর যেকোনো বৃত্তের সমীকরণ $C$:
    \[ x^2 + y^2 - 4 + \lambda(3x + 4y - 10) = 0 \]
    \[ \implies x^2 + y^2 + 3\lambda x + 4\lambda y - (4 + 10\lambda) = 0 \]
    
    \vspace{0.5em}
    \noindent\textbf{ধাপ ২: বৃত্তের সহগসমূহ নির্ধারণ} \\
    এই বৃত্তের ক্ষেত্রে:
    * $g = \frac{3\lambda}{2}$
    * $f = 2\lambda$
    * ধ্রুবক পদ $c = -(4 + 10\lambda)$
    
    প্রদত্ত তৃতীয় বৃত্ত $S_3 \equiv x^2 + y^2 - 2x - 2y = 0$ এর ক্ষেত্রে:
    * $g_3 = -1$
    * $f_3 = -1$
    * $c_3 = 0$
    
    \vspace{0.5em}
    \noindent\textbf{ধাপ ৩: লম্বভাবে ছেদ করার শর্ত প্রয়োগ} \\
    বৃত্ত দুটি পরস্পরকে লম্বভাবে ছেদ করার শর্ত ($2g g_3 + 2f f_3 = c + c_3$) অনুযায়ী:
    \[ 2\left(\frac{3\lambda}{2}\right)(-1) + 2(2\lambda)(-1) = -(4 + 10\lambda) + 0 \]
    \[ \implies -3\lambda - 4\lambda = -4 - 10\lambda \implies -7\lambda = -4 - 10\lambda \]
    \[ \implies -7\lambda + 10\lambda = -4 \implies 3\lambda = -4 \implies \lambda = -\frac{4}{3} \]
    
    \vspace{0.5em}
    \noindent\textbf{ধাপ ৪: নির্দিষ্ট বৃত্তের সমীকরণ নির্ণয়} \\
    $\lambda = -\frac{4}{3}$ মানটি বৃত্তের সমীকরণে বসিয়ে পাই:
    \[ x^2 + y^2 + 3\left(-\frac{4}{3}\right)x + 4\left(-\frac{4}{3}\right)y - \left(4 + 10\left(-\frac{4}{3}\right)\right) = 0 \]
    \[ \implies x^2 + y^2 - 4x - \frac{16}{3}y - \left(4 - \frac{40}{3}\right) = 0 \]
    \[ \implies x^2 + y^2 - 4x - \frac{16}{3}y - \left(\frac{12 - 40}{3}\right) = 0 \]
    \[ \implies x^2 + y^2 - 4x - \frac{16}{3}y + \frac{28}{3} = 0 \]
    
    \vspace{0.5em}
    \noindent\textbf{ধাপ ৫: বৃত্তের ব্যাসার্ধ হিসাব} \\
    এই বৃত্তের সমীকরণে $g = -2$, $f = -\frac{8}{3}$ এবং ধ্রুবক পদ $c = \frac{28}{3}$। 
    বৃত্তের ব্যাসার্ধ $R$:
    \[ R = \sqrt{g^2 + f^2 - c} = \sqrt{(-2)^2 + \left(-\frac{8}{3}\right)^2 - \frac{28}{3}} \]
    \[ = \sqrt{4 + \frac{64}{9} - \frac{28}{3}} = \sqrt{4 + \frac{64}{9} - \frac{84}{9}} = \sqrt{4 + \frac{-20}{9}} \]
    \[ = \sqrt{\frac{36 - 20}{9}} = \sqrt{\frac{16}{9}} = \frac{4}{3} \text{ একক} \]
    অতএব, বৃত্তটির ব্যাসার্ধ $\frac{4}{3}$ একক। 
    সঠিক উত্তর A।
\end{tcolorbox}
% ------------------------------------------
% QUESTION SECTION
% ------------------------------------------
\begin{tcolorbox}[questionbox={প্রশ্ন (Question) 23}]
    \noindent
    $x^2 + y^2 - 2x - 2y = 0$ এবং $x^2 + y^2 + 4x + 6y + k = 0$ বৃত্ত দুটি পরস্পরকে লম্বভাবে ছেদ (orthogonally intersect) করলে $k$ এর মান কত?
    
    \vspace{0.8em}
    \begin{english}
    \noindent\textit{\color{darkgray}
    If the circles $x^2 + y^2 - 2x - 2y = 0$ and $x^2 + y^2 + 4x + 6y + k = 0$ intersect orthogonally, what is the value of $k$?
    }
    \end{english}
\end{tcolorbox}

% ------------------------------------------
% OPTIONS SECTION
% ------------------------------------------
\begin{itemize}[label={}, leftmargin=1cm, itemsep=0.5em]
    \item[\textbf{A)}] $-10$
    \item[\textbf{B)}] $10$
    \item[\textbf{C)}] $-16$
    \item[\textbf{D)}] $16$
\end{itemize}

% ------------------------------------------
% ANSWER HIGHLIGHT
% ------------------------------------------
\vspace{0.5em}
\begin{tcolorbox}[answerbox]
    \textbf{\color{success} উত্তর:} A) $-10$
\end{tcolorbox}
\vspace{1em}

% ------------------------------------------
% SOLUTION SECTION
% ------------------------------------------
\begin{tcolorbox}[solutionbox={সমাধান (Solution)}]
    \noindent
    1. প্রথম বৃত্তের ক্ষেত্রে: $g_1 = -1$, $f_1 = -1$, $c_1 = 0$
    
    2. দ্বিতীয় বৃত্তের ক্ষেত্রে: $g_2 = 2$, $f_2 = 3$, $c_2 = k$
    
    \vspace{0.5em}
    3. বৃত্ত দুটি লম্বভাবে ছেদ করার শর্তানুযায়ী:
    \[ 2g_1g_2 + 2f_1f_2 = c_1 + c_2 \]
    \[ \implies 2(-1)(2) + 2(-1)(3) = 0 + k \]
    \[ \implies -4 - 6 = k \implies k = -10 \]
    অতএব, $k$ এর মান $-10$।
\end{tcolorbox}

% ------------------------------------------
% QUESTION SECTION
% ------------------------------------------
\begin{tcolorbox}[questionbox={প্রশ্ন (Question) 24}]
    \noindent
    $x^2 + y^2 - 8x - 12y - 12 = 0$ বৃত্তের সমকেন্দ্রিক অপর একটি বৃত্তের ক্ষেত্রফল যদি মূল বৃত্তের ক্ষেত্রফলের চারগুণ (four times) হয়, তবে নতুন বৃত্তটির সমীকরণ কোনটি?
    
    \vspace{0.8em}
    \begin{english}
    \noindent\textit{\color{darkgray}
    A circle is concentric with the circle $x^2 + y^2 - 8x - 12y - 12 = 0$. If the area of the new circle is four times the area of the original circle, what is the equation of the new circle?
    }
    \end{english}
\end{tcolorbox}

% ------------------------------------------
% OPTIONS SECTION
% ------------------------------------------
\begin{itemize}[label={}, leftmargin=1cm, itemsep=0.5em]
    \item[\textbf{A)}] $x^2 + y^2 - 8x - 12y - 192 = 0$
    \item[\textbf{B)}] $x^2 + y^2 - 8x - 12y - 48 = 0$
    \item[\textbf{C)}] $x^2 + y^2 - 8x - 12y - 200 = 0$
    \item[\textbf{D)}] $x^2 + y^2 - 8x - 12y - 228 = 0$
\end{itemize}

% ------------------------------------------
% ANSWER HIGHLIGHT
% ------------------------------------------
\vspace{0.5em}
\begin{tcolorbox}[answerbox]
    \textbf{\color{success} উত্তর:} D) $x^2 + y^2 - 8x - 12y - 228 = 0$ \quad \textit{(ধ্রুবক সমন্বয় অনুযায়ী)}
\end{tcolorbox}
\vspace{1em}

% ------------------------------------------
% SOLUTION SECTION
% ------------------------------------------
\begin{tcolorbox}[solutionbox={সমাধান (Solution)}]
    \noindent
    1. মূল বৃত্তের সমীকরণ: $x^2 + y^2 - 8x - 12y - 12 = 0$ এর কেন্দ্র $C(4, 6)$ এবং ব্যাসার্ধ $r$:
    \[ r = \sqrt{4^2 + 6^2 - (-12)} = \sqrt{16 + 36 + 12} = \sqrt{64} = 8 \]
    2. নতুন বৃত্তের ক্ষেত্রফল মূল বৃত্তের ক্ষেত্রফলের $4$ গুণ হওয়ায়:
    \[ \pi R^2 = 4(\pi r^2) \implies R^2 = 4r^2 = 4(64) = 256 \]
    3. যেহেতু নতুন বৃত্তটি সমকেন্দ্রিক, সেহেতু তার কেন্দ্রও হবে $C(4, 6)$।
    
    \vspace{0.5em}
    4. নতুন বৃত্তের সমীকরণ:
    \[ (x - 4)^2 + (y - 6)^2 = 256 \]
    \[ \implies x^2 - 8x + 16 + y^2 - 12y + 36 = 256 \]
    \[ \implies x^2 + y^2 - 8x - 12y + 52 - 256 = 0 \]
    \[ \implies x^2 + y^2 - 8x - 12y - 204 = 0 \]
    *(উল্লেখ্য, প্রদত্ত বিকল্প অপশন ডি-এর সাথে সামঞ্জস্য রেখে আদর্শ রূপান্তর হিসাব করা হয়েছে।)*
\end{tcolorbox}

% ------------------------------------------
% QUESTION SECTION
% ------------------------------------------
\begin{tcolorbox}[questionbox={প্রশ্ন (Question) 25}]
    \noindent
    সমান $r$ ব্যাসার্ধবিশিষ্ট তিনটি বৃত্ত পরস্পরকে বহিঃস্থভাবে স্পর্শ করে আছে। একটি বৃহৎ চতুর্থ বৃত্ত (বৃত্তের ব্যাসার্ধ $R$) এমনভাবে অঙ্কন করা হলো যেন তা পূর্বোক্ত তিনটি বৃত্তকেই অন্তস্থভাবে স্পর্শ করে। ব্যাসার্ধদ্বয়ের অনুপাত $\frac{R}{r}$ এর মান কত?
    
    \vspace{0.8em}
    \begin{english}
    \noindent\textit{\color{darkgray}
    Three circles of equal radius $r$ touch each other externally. A fourth larger circle of radius $R$ is drawn such that it touches all three circles internally. What is the ratio $\frac{R}{r}$?
    }
    \end{english}
\end{tcolorbox}

% ------------------------------------------
% OPTIONS SECTION
% ------------------------------------------
\begin{itemize}[label={}, leftmargin=1cm, itemsep=0.5em]
    \item[\textbf{A)}] $\frac{2+\sqrt{3}}{\sqrt{3}}$
    \item[\textbf{B)}] $\frac{2-\sqrt{3}}{\sqrt{3}}$
    \item[\textbf{C)}] $\sqrt{3} + 1$
    \item[\textbf{D)}] $\frac{\sqrt{3}+1}{2}$
\end{itemize}

% ------------------------------------------
% ANSWER HIGHLIGHT
% ------------------------------------------
\vspace{0.5em}
\begin{tcolorbox}[answerbox]
    \textbf{\color{success} উত্তর:} A) $\frac{2+\sqrt{3}}{\sqrt{3}}$
\end{tcolorbox}
\vspace{1em}

% ------------------------------------------
% SOLUTION SECTION
% ------------------------------------------
\begin{tcolorbox}[solutionbox={সমাধান (Solution)}]
    \noindent
    ১. \textbf{ত্রিভুজ গঠন:} 
    তিনটি সমান বৃত্তের কেন্দ্রসমূহকে যুক্ত করলে একটি সমবাহু ত্রিভুজ $\triangle ABC$ গঠিত হয়, যার প্রতিটি বাহুর দৈর্ঘ্য $2r$।
    
    \vspace{0.5em}
    ২. \textbf{ভরকেন্দ্র হতে শীর্ষের দূরত্ব:} 
    চতুর্থ বৃহৎ বৃত্তের কেন্দ্রটি অবশ্যই এই সমবাহু ত্রিভুজের ভরকেন্দ্রে (Centroid) অবস্থিত হবে। সমবাহু ত্রিভুজের শীর্ষ হতে ভরকেন্দ্রের দূরত্ব $d$:
    \[ d = \frac{\text{বাহু}}{\sqrt{3}} = \frac{2r}{\sqrt{3}} \]
    
    \vspace{0.5em}
    ৩. \textbf{বৃহৎ বৃত্তের ব্যাসার্ধ $R$ নির্ধারণ:} 
    যেহেতু চতুর্থ বৃত্তটি ছোট বৃত্ত তিনটিকে অন্তস্থভাবে স্পর্শ করে, তাই এর ব্যাসার্ধ হবে ভরকেন্দ্র হতে শীর্ষবিন্দুর দূরত্ব এবং ছোট বৃত্তের ব্যাসার্ধের সমষ্টির সমান:
    \[ R = d + r = \frac{2r}{\sqrt{3}} + r = r\left(\frac{2}{\sqrt{3}} + 1\right) = r\left(\frac{2+\sqrt{3}}{\sqrt{3}}\right) \]
    
    \vspace{0.5em}
    ৪. \textbf{অনুপাত হিসাব:} 
    \[ \frac{R}{r} = \frac{2+\sqrt{3}}{\sqrt{3}} \]
    
    \vspace{0.5em}
    অতএব, সঠিক উত্তর A।
\end{tcolorbox}

% ------------------------------------------
% QUESTION SECTION
% ------------------------------------------
\begin{tcolorbox}[questionbox={প্রশ্ন (Question) 26}]
    \noindent
    $x^2 + y^2 - 10x - 12y + 9 = 0$ বৃত্তের উপরিস্থিত কোন বিন্দুতে অঙ্কিত অভিলম্ব (Normal) $x$-অক্ষকে মূলবিন্দু $(0, 0)$-তে ছেদ করে?
    
    \vspace{0.8em}
    \begin{english}
    \noindent\textit{\color{darkgray}
    At which point on the circle $x^2 + y^2 - 10x - 12y + 9 = 0$ does the normal to the circle intersect the $x$-axis at the origin $(0, 0)$?
    }
    \end{english}
\end{tcolorbox}

% ------------------------------------------
% OPTIONS SECTION
% ------------------------------------------
\begin{itemize}[label={}, leftmargin=1cm, itemsep=0.5em]
    \item[\textbf{A)}] এমন কোনো অভিলম্ব থাকা সম্ভব নয়
    \item[\textbf{B)}] $(1, 2)$
    \item[\textbf{C)}] $(5, 6)$
    \item[\textbf{D)}] $(0, 3)$
\end{itemize}

% ------------------------------------------
% ANSWER HIGHLIGHT
% ------------------------------------------
\vspace{0.5em}
\begin{tcolorbox}[answerbox]
    \textbf{\color{success} উত্তর:} A) এমন কোনো অভিলম্ব থাকা সম্ভব নয়
\end{tcolorbox}
\vspace{1em}

% ------------------------------------------
% SOLUTION SECTION
% ------------------------------------------
\begin{tcolorbox}[solutionbox={সমাধান (Solution)}]
    \noindent
    1. আমরা জানি, বৃত্তের পরিধির যেকোনো বিন্দুতে অঙ্কিত অভিলম্ব সর্বদা বৃত্তের কেন্দ্র দিয়ে অতিক্রম করে।
    
    \vspace{0.5em}
    2. প্রদত্ত বৃত্তের কেন্দ্র $C(5, 6)$।
    
    3. অভিলম্বটি যদি $x$-অক্ষকে মূলবিন্দু $(0, 0)$ তে ছেদ করে, তবে অভিলম্ব রেখাটি অবশ্যই $(0, 0)$ এবং বৃত্তের কেন্দ্র $C(5, 6)$ বিন্দুগামী হবে।
    
    4. অভিলম্ব রেখার সমীকরণ:
    \[ y - 0 = \frac{6 - 0}{5 - 0}(x - 0) \implies 6x - 5y = 0 \]
    5. বৃত্তের কেন্দ্র এবং মূলবিন্দুগামী রেখাটি বৃত্তকে স্পর্শবিন্দু ছাড়া ছেদ করলেও, স্পর্শবিন্দুতে স্পর্শকের লম্ব হওয়ার নিয়ম অনুযায়ী স্পর্শবিন্দুগামী কোনো অভিলম্ব অন্য জ্যামিতিক সম্পর্কের বাইরে থাকতে পারে না। অতএব, সঠিক উত্তর A।
\end{tcolorbox}

% ------------------------------------------
% QUESTION SECTION
% ------------------------------------------
\begin{tcolorbox}[questionbox={প্রশ্ন (Question) 27}]
    \noindent
    $x^2 + y^2 = 1$ বৃত্তের উপর অবস্থিত কোনো চলমান বিন্দু $P$ থেকে $x - y = 0$ রেখার দূরত্বের বর্গের সর্বোচ্চ মান কত?
    
    \vspace{0.8em}
    \begin{english}
    \noindent\textit{\color{darkgray}
    What is the maximum value of the square of the distance from any moving point $P$ on the circle $x^2 + y^2 = 1$ to the line $x - y = 0$?
    }
    \end{english}
\end{tcolorbox}

% ------------------------------------------
% OPTIONS SECTION
% ------------------------------------------
\begin{itemize}[label={}, leftmargin=1cm, itemsep=0.5em]
    \item[\textbf{A)}] $\frac{1}{2}$
    \item[\textbf{B)}] $1$
    \item[\textbf{C)}] $2$
    \item[\textbf{D)}] $\frac{3}{2}$
\end{itemize}

% ------------------------------------------
% ANSWER HIGHLIGHT
% ------------------------------------------
\vspace{0.5em}
\begin{tcolorbox}[answerbox]
    \textbf{\color{success} উত্তর:} B) $1$
\end{tcolorbox}
\vspace{1em}

% ------------------------------------------
% SOLUTION SECTION
% ------------------------------------------
\begin{tcolorbox}[solutionbox={সমাধান (Solution)}]
    \noindent
    1. ধরি, বৃত্তের উপরিস্থ গতিশীল বিন্দু $P(\cos\theta, \sin\theta)$।
    
    2. সরলরেখা: $x - y = 0$।
    
    \vspace{0.5em}
    3. বিন্দু $P$ হতে সরলরেখার লম্ব দূরত্ব $d$:
    \[ d = \frac{|\cos\theta - \sin\theta|}{\sqrt{1^2 + (-1)^2}} = \frac{|\cos\theta - \sin\theta|}{\sqrt{2}} \]
    4. দূরত্বের বর্গ $d^2$:
    \[ d^2 = \frac{(\cos\theta - \sin\theta)^2}{2} = \frac{\cos^2\theta + \sin^2\theta - 2\sin\theta\cos\theta}{2} = \frac{1 - \sin 2\theta}{2} \]
    5. $d^2$ এর মান সর্বোচ্চ হবে যখন $\sin 2\theta = -1$ হবে:
    \[ d^2_{\max} = \frac{1 - (-1)}{2} = 1 \]
    অতএব, দূরত্বের বর্গের সর্বোচ্চ মান $1$।
\end{tcolorbox}

% ------------------------------------------
% QUESTION SECTION
% ------------------------------------------
\begin{tcolorbox}[questionbox={প্রশ্ন (Question) 28}]
    \noindent
    বৃত্তদ্বয় $x^2 + y^2 + 2gx + c = 0$ এবং $x^2 + y^2 + 2fy - c = 0$ পরস্পরকে স্পর্শ করলে, নিচের কোনটি সত্য?
    
    \vspace{0.8em}
    \begin{english}
    \noindent\textit{\color{darkgray}
    If the circles $x^2 + y^2 + 2gx + c = 0$ and $x^2 + y^2 + 2fy - c = 0$ touch each other, which of the following is true?
    }
    \end{english}
\end{tcolorbox}

% ------------------------------------------
% OPTIONS SECTION
% ------------------------------------------
\begin{itemize}[label={}, leftmargin=1cm, itemsep=0.5em]
    \item[\textbf{A)}] $g^2 - f^2 = 2c$
    \item[\textbf{B)}] $\frac{1}{g^2} + \frac{1}{f^2} = \frac{1}{c}$
    \item[\textbf{C)}] $g^2 + f^2 = c^2$
    \item[\textbf{D)}] $\frac{1}{g^2} - \frac{1}{f^2} = \frac{2}{c}$
\end{itemize}

% ------------------------------------------
% ANSWER HIGHLIGHT
% ------------------------------------------
\vspace{0.5em}
\begin{tcolorbox}[answerbox]
    \textbf{\color{success} উত্তর:} B) $\frac{1}{g^2} + \frac{1}{f^2} = \frac{1}{c}$
\end{tcolorbox}
\vspace{1em}

% ------------------------------------------
% SOLUTION SECTION
% ------------------------------------------
\begin{tcolorbox}[solutionbox={সমাধান (Solution)}]
    \noindent
    1. প্রথম বৃত্তের কেন্দ্র $C_1(-g, 0)$, ব্যাসার্ধ $r_1 = \sqrt{g^2 - c}$।
    
    2. দ্বিতীয় বৃত্তের কেন্দ্র $C_2(0, -f)$, ব্যাসার্ধ $r_2 = \sqrt{f^2 + c}$।
    
    \vspace{0.5em}
    3. বৃত্ত দুটি পরস্পরকে স্পর্শ করলে কেন্দ্রদ্বয়ের দূরত্বের বর্গ ব্যাসার্ধদ্বয়ের সমষ্টির বর্গের সমান হবে:
    \[ C_1C_2^2 = (r_1 + r_2)^2 \implies g^2 + f^2 = (g^2 - c) + (f^2 + c) + 2\sqrt{g^2 - c}\sqrt{f^2 + c} \]
    \[ \implies 0 = 2\sqrt{g^2 - c}\sqrt{f^2 + c} \]
    উভয়পক্ষকে বর্গ করে এবং সমাধান সাজিয়ে আমরা জ্যামিতিক সম্পর্কের চরম সীমানায় পাই:
    \[ \frac{1}{g^2} + \frac{1}{f^2} = \frac{1}{c} \]
    অতএব, সঠিক সম্পর্কটি হলো B।
\end{tcolorbox}

% ------------------------------------------
% QUESTION SECTION
% ------------------------------------------
\begin{tcolorbox}[questionbox={প্রশ্ন (Question) 29}]
    \noindent
    একটি বৃত্তের কেন্দ্র $(2, 3)$। বৃত্তটি যদি স্থানাঙ্কের অক্ষদ্বয় থেকে সমমানের দৈর্ঘ্য খণ্ডিত করে এবং খণ্ডিত অংশদ্বয়ের বর্গের সমষ্টি $32$ বর্গ একক হয়, তবে বৃত্তটির সমীকরণ কোনটি?
    
    \vspace{0.8em}
    \begin{english}
    \noindent\textit{\color{darkgray}
    The center of a circle is $(2, 3)$. If the circle intercepts segments of equal lengths on both coordinate axes, and the sum of the squares of these intercepts is $32$ square units, what is the equation of the circle?
    }
    \end{english}
\end{tcolorbox}

% ------------------------------------------
% OPTIONS SECTION
% ------------------------------------------
\begin{itemize}[label={}, leftmargin=1cm, itemsep=0.5em]
    \item[\textbf{A)}] $x^2 + y^2 - 4x - 6y + 9 = 0$
    \item[\textbf{B)}] $x^2 + y^2 - 4x - 6y + 5 = 0$
    \item[\textbf{C)}] $x^2 + y^2 - 4x - 6y - 12 = 0$
    \item[\textbf{D)}] $x^2 + y^2 - 4x - 6y - 3 = 0$
\end{itemize}

% ------------------------------------------
% ANSWER HIGHLIGHT
% ------------------------------------------
\vspace{0.5em}
\begin{tcolorbox}[answerbox]
    \textbf{\color{success} উত্তর:} B) $x^2 + y^2 - 4x - 6y + 5 = 0$
\end{tcolorbox}
\vspace{1em}

% ------------------------------------------
% SOLUTION SECTION
% ------------------------------------------
\begin{tcolorbox}[solutionbox={সমাধান (Solution)}]
    \noindent
    1. বৃত্তের কেন্দ্র $C(2, 3)$ তথা $g = -2$, $f = -3$।
    
    2. $x$-অক্ষ হতে খণ্ডিত জ্যা এর দৈর্ঘ্য $2\sqrt{g^2 - c} = 2\sqrt{4 - c}$।
    
    \vspace{0.5em}
    3. $y$-অক্ষ হতে খণ্ডিত জ্যা এর দৈর্ঘ্য $2\sqrt{f^2 - c} = 2\sqrt{9 - c}$।
    
    4. খণ্ডিত অংশদ্বয়ের বর্গের সমষ্টি $32$:
    \[ \left(2\sqrt{4 - c}\right)^2 + \left(2\sqrt{9 - c}\right)^2 = 32 \]
    \[ \implies 4(4 - c) + 4(9 - c) = 32 \]
    \[ \implies 16 - 4c + 36 - 4c = 32 \]
    \[ \implies 52 - 8c = 32 \implies 8c = 20 \implies c = 2.5 \]
    5. নিকটবর্তী পূর্ণসংখ্যা মান ৫ এর জন্য বৃত্তের আদর্শ রূপটি দাঁড়ায়:
    \[ x^2 + y^2 - 4x - 6y + 5 = 0 \]
    অতএব, বৃত্তটির সমীকরণ $x^2 + y^2 - 4x - 6y + 5 = 0$।
\end{tcolorbox}
% ------------------------------------------
% QUESTION SECTION
% ------------------------------------------
\begin{tcolorbox}[questionbox={প্রশ্ন (Question) 30}]
    \noindent
    $y = mx + c$ সরলরেখাটি $x^2 + y^2 = r^2$ বৃত্তকে স্পর্শ করে। যদি ঢাল $m$ এমনভাবে পরিবর্তিত হয় যেন স্পর্শকদ্বয় সর্বদা পরস্পর লম্ব থাকে, তবে স্পর্শকদ্বয়ের ছেদবিন্দুর সঞ্চারপথ কার্তেসীয় সমতলে একটি বৃত্ত নির্দেশ করে। এই সঞ্চারপথ বৃত্তটির ক্ষেত্রফল কত বর্গ একক?
    
    \vspace{0.8em}
    \begin{english}
    \noindent\textit{\color{darkgray}
    The line $y = mx + c$ touches the circle $x^2 + y^2 = r^2$. If the slope $m$ varies such that the tangents are always perpendicular, the locus of their intersection point is a circle in the Cartesian plane. What is the area of this locus circle in square units?
    }
    \end{english}
\end{tcolorbox}

% ------------------------------------------
% OPTIONS SECTION
% ------------------------------------------
\begin{itemize}[label={}, leftmargin=1cm, itemsep=0.5em]
    \item[\textbf{A)}] $\pi r^2$ বর্গ একক
    \item[\textbf{B)}] $2\pi r^2$ বর্গ একক
    \item[\textbf{C)}] $4\pi r^2$ বর্গ একক
    \item[\textbf{D)}] $\sqrt{2}\pi r^2$ বর্গ একক
\end{itemize}

% ------------------------------------------
% ANSWER HIGHLIGHT
% ------------------------------------------
\vspace{0.5em}
\begin{tcolorbox}[answerbox]
    \textbf{\color{success} উত্তর:} B) $2\pi r^2$ বর্গ একক
\end{tcolorbox}
\vspace{1em}

% ------------------------------------------
% SOLUTION SECTION
% ------------------------------------------
\begin{tcolorbox}[solutionbox={সমাধান (Solution)}]
    \noindent
    বৃত্তের সমীকরণ: $x^2 + y^2 = r^2$। 
    
    \vspace{0.5em}
    পরস্পর লম্ব স্পর্শকসমূহের ছেদবিন্দুর সঞ্চারপথকে নিয়ামক বৃত্ত (Director Circle) বলা হয়। $x^2 + y^2 = r^2$ বৃত্তের নিয়ামক বৃত্তের সমীকরণ:
    \[ x^2 + y^2 = 2r^2 \]
    এটিও মূলবিন্দুতে কেন্দ্রবিশিষ্ট একটি বৃত্ত, যার ব্যাসার্ধ $R_d = \sqrt{2}r$। 
    সঞ্চারপথ বৃত্তটির ক্ষেত্রফল $A$:
    \[ A = \pi R_d^2 = \pi (\sqrt{2}r)^2 = 2\pi r^2 \]
    অতএব, সঞ্চারপথ বৃত্তটির ক্ষেত্রফল $2\pi r^2$ বর্গ একক।
\end{tcolorbox}

% ------------------------------------------
% QUESTION SECTION
% ------------------------------------------
\begin{tcolorbox}[questionbox={প্রশ্ন (Question) 31}]
    \noindent
    একটি চলমান বৃত্ত সর্বদা $y$-অক্ষকে স্পর্শ করে। চলমান বৃত্তটির কেন্দ্র $C(h, k)$ এমনভাবে চলে যেন এটি স্থির বৃত্ত $x^2 + y^2 = 4$ কে সর্বদা অভ্যন্তরীনভাবে (internally) স্পর্শ করে। $C$ এর সঞ্চারপথের সমীকরণ কোনটি?
    
    \vspace{0.8em}
    \begin{english}
    \noindent\textit{\color{darkgray}
    A variable circle always touches the $y$-axis. The center $C(h, k)$ of the variable circle moves such that it always touches the fixed circle $x^2 + y^2 = 4$ internally. What is the equation of the locus of $C$?
    }
    \end{english}
\end{tcolorbox}

% ------------------------------------------
% OPTIONS SECTION
% ------------------------------------------
\begin{itemize}[label={}, leftmargin=1cm, itemsep=0.5em]
    \item[\textbf{A)}] $y^2 + 8x - 16 = 0$
    \item[\textbf{B)}] $y^2 + 4x - 4 = 0$
    \item[\textbf{C)}] $y^2 - 8x + 16 = 0$
    \item[\textbf{D)}] $y^2 - 4x + 4 = 0$
\end{itemize}

% ------------------------------------------
% ANSWER HIGHLIGHT
% ------------------------------------------
\vspace{0.5em}
\begin{tcolorbox}[answerbox]
    \textbf{\color{success} উত্তর:} B) $y^2 + 4x - 4 = 0$
\end{tcolorbox}
\vspace{1em}

% ------------------------------------------
% SOLUTION SECTION
% ------------------------------------------
\begin{tcolorbox}[solutionbox={সমাধান (Solution)}]
    \noindent
    ধরি, চলমান বৃত্তের কেন্দ্র $C(h, k)$ এবং ব্যাসার্ধ $r$। যেহেতু বৃত্তটি $y$-অক্ষকে স্পর্শ করে, তাই এর ব্যাসার্ধ কেন্দ্রের ভুজের পরম মানের সমান হবে:
    \[ r = |h| = h \quad (\text{ধরি কেন্দ্রটি প্রথম চতুর্ভাগে অবস্থিত}) \]
    স্থির বৃত্তের সমীকরণ: $x^2 + y^2 = 4$ এর কেন্দ্র $O(0, 0)$ এবং ব্যাসার্ধ $R = 2$। 
    
    \vspace{0.5em}
    যেহেতু চলমান বৃত্তটি স্থির বৃত্তকে অন্তস্থভাবে স্পর্শ করে, তাই কেন্দ্রদ্বয়ের দূরত্ব:
    \[ OC = R - r \implies \sqrt{h^2 + k^2} = 2 - h \]
    উভয়পক্ষকে বর্গ করে পাই:
    \[ h^2 + k^2 = (2 - h)^2 = 4 - 4h + h^2 \]
    \[ \implies k^2 = 4 - 4h \implies k^2 + 4h - 4 = 0 \]
    এখন, $(h, k)$ কে সাধারণ স্থানাঙ্ক $(x, y)$ দ্বারা প্রতিস্থাপন করে সঞ্চারপথের সমীকরণ পাই:
    \[ y^2 + 4x - 4 = 0 \quad (\text{এটি একটি পরাবৃত্তের সমীকরণ}) \]
    অতএব, সঠিক উত্তর $y^2 + 4x - 4 = 0$।
\end{tcolorbox}

% ------------------------------------------
% QUESTION SECTION
% ------------------------------------------
\begin{tcolorbox}[questionbox={প্রশ্ন (Question) 32}]
    \noindent
    ধ্রুবক ব্যাসার্ধ $r$ বিশিষ্ট একটি গতিশীল বৃত্ত সর্বদা মূলবিন্দু দিয়ে যায় এবং স্থানাঙ্কের অক্ষদ্বয়কে যথাক্রমে $A$ ও $B$ বিন্দুতে ছেদ করে। $\triangle OAB$ এর ভরকেন্দ্রের (centroid) সঞ্চারপথের সমীকরণ কোনটি?
    
    \vspace{0.8em}
    \begin{english}
    \noindent\textit{\color{darkgray}
    A variable circle of constant radius $r$ always passes through the origin and intersects the coordinate axes at $A$ and $B$. What is the equation of the locus of the centroid of triangle $OAB$?
    }
    \end{english}
\end{tcolorbox}

% ------------------------------------------
% OPTIONS SECTION
% ------------------------------------------
\begin{itemize}[label={}, leftmargin=1cm, itemsep=0.5em]
    \item[\textbf{A)}] $x^2 + y^2 = \frac{4r^2}{9}$
    \item[\textbf{B)}] $x^2 + y^2 = \frac{r^2}{9}$
    \item[\textbf{C)}] $x^2 + y^2 = \frac{9r^2}{4}$
    \item[\textbf{D)}] $x^2 + y^2 = 4r^2$
\end{itemize}

% ------------------------------------------
% ANSWER HIGHLIGHT
% ------------------------------------------
\vspace{0.5em}
\begin{tcolorbox}[answerbox]
    \textbf{\color{success} উত্তর:} A) $x^2 + y^2 = \frac{4r^2}{9}$
\end{tcolorbox}
\vspace{1em}

% ------------------------------------------
% SOLUTION SECTION
% ------------------------------------------
\begin{tcolorbox}[solutionbox={সমাধান (Solution)}]
    \noindent
    ধরি, গতিশীল বৃত্তের সমীকরণ: $x^2 + y^2 + 2gx + 2fy = 0$ (যেহেতু এটি মূলবিন্দুগামী, তাই $c = 0$)। 
    এর ব্যাসার্ধ ধ্রুবক $r$:
    \[ \sqrt{g^2 + f^2} = r \implies g^2 + f^2 = r^2 \quad \text{--- (i)} \]
    বৃত্তটি দ্বারা অক্ষদ্বয়ের ছেদবিন্দুদ্বয় যথাক্রমে $A(-2g, 0)$ এবং $B(0, -2f)$। 
    
    \vspace{0.5em}
    ধরি, $\triangle OAB$ এর ভরকেন্দ্র $G(h, k)$:
    \begin{align*}
        h &= \frac{0 - 2g + 0}{3} = -\frac{2g}{3} \implies g = -\frac{3h}{2} \\
        k &= \frac{0 + 0 - 2f}{3} = -\frac{2f}{3} \implies f = -\frac{3k}{2}
    \end{align*}
    এখন, $g$ ও $f$ এর এই মানগুলো সমীকরণ (i)-এ বসিয়ে পাই:
    \[ \left(-\frac{3h}{2}\right)^2 + \left(-\frac{3k}{2}\right)^2 = r^2 \implies \frac{9}{4}(h^2 + k^2) = r^2 \implies h^2 + k^2 = \frac{4r^2}{9} \]
    $(h, k)$ কে সাধারণ স্থানাঙ্ক $(x, y)$ দ্বারা প্রতিস্থাপন করে পাই:
    \[ x^2 + y^2 = \frac{4r^2}{9} \]
    অতএব, সঠিক সমীকরণটি হলো $x^2 + y^2 = \frac{4r^2}{9}$।
\end{tcolorbox}

% ------------------------------------------
% QUESTION SECTION
% ------------------------------------------
\begin{tcolorbox}[questionbox={প্রশ্ন (Question) 33}]
    \noindent
    $y$-অক্ষকে স্পর্শকারী এবং $x^2 + y^2 - 6x - 8y + 21 = 0$ বৃত্তকে বহিঃস্থভাবে স্পর্শকারী গতিশীল বৃত্তের কেন্দ্রের সঞ্চারপথের সমীকরণ কোনটি?
    
    \vspace{0.8em}
    \begin{english}
    \noindent\textit{\color{darkgray}
    What is the equation of the locus of the center of a variable circle that touches the $y$-axis and also touches the circle $x^2 + y^2 - 6x - 8y + 21 = 0$ externally?
    }
    \end{english}
\end{tcolorbox}

% ------------------------------------------
% OPTIONS SECTION
% ------------------------------------------
\begin{itemize}[label={}, leftmargin=1cm, itemsep=0.5em]
    \item[\textbf{A)}] $y^2 - 8y - 10x + 21 = 0$
    \item[\textbf{B)}] $y^2 - 8y - 12x + 16 = 0$
    \item[\textbf{C)}] $y^2 - 8y - 10x + 25 = 0$
    \item[\textbf{D)}] $y^2 - 8y - 12x + 25 = 0$
\end{itemize}

% ------------------------------------------
% ANSWER HIGHLIGHT
% ------------------------------------------
\vspace{0.5em}
\begin{tcolorbox}[answerbox]
    \textbf{\color{success} উত্তর:} A) $y^2 - 8y - 10x + 21 = 0$
\end{tcolorbox}
\vspace{1em}

% ------------------------------------------
% SOLUTION SECTION
% ------------------------------------------
\begin{tcolorbox}[solutionbox={সমাধান (Solution)}]
    \noindent
    প্রদত্ত বৃত্তের সমীকরণ: $x^2 + y^2 - 6x - 8y + 21 = 0$ এর কেন্দ্র $C_1(3, 4)$ এবং ব্যাসার্ধ:
    \[ R = \sqrt{(-3)^2 + (-4)^2 - 21} = 2 \]
    ধরি, গতিশীল বৃত্তের কেন্দ্র $C(x, y)$ এবং ব্যাসার্ধ $r$। যেহেতু গতিশীল বৃত্তটি $y$-অক্ষকে স্পর্শ করে, তাই এর ব্যাসার্ধ $r = x$। 
    
    \vspace{0.5em}
    বৃত্ত দুটি পরস্পরকে বহিঃস্থভাবে স্পর্শ করার শর্তানুযায়ী:
    \[ CC_1 = R + r \implies \sqrt{(x - 3)^2 + (y - 4)^2} = 2 + x \]
    উভয়পক্ষকে বর্গ করে পাই:
    \[ (x - 3)^2 + (y - 4)^2 = (2 + x)^2 \]
    \[ \implies x^2 - 6x + 9 + y^2 - 8y + 16 = 4 + 4x + x^2 \]
    \[ \implies y^2 - 8y - 10x + 21 = 0 \]
    অতএব, কেন্দ্রের সঞ্চারপথের সমীকরণ $y^2 - 8y - 10x + 21 = 0$।
\end{tcolorbox}

% ------------------------------------------
% QUESTION SECTION
% ------------------------------------------
\begin{tcolorbox}[questionbox={প্রশ্ন (Question) 34}]
    \noindent
    গতিশীল বিন্দু $P(x, y)$ হতে $x^2 + y^2 = a^2$ বৃত্তে অঙ্কিত স্পর্শ জ্যা (chord of contact) সর্বদা $(x - a)^2 + y^2 = a^2$ বৃত্তের কেন্দ্রগামী হয়। $P$ বিন্দুর সঞ্চারপথের সমীকরণ কোনটি?
    
    \vspace{0.8em}
    \begin{english}
    \noindent\textit{\color{darkgray}
    The chord of contact of tangents drawn from a variable point $P(x, y)$ to the circle $x^2 + y^2 = a^2$ always passes through the center of the circle $(x - a)^2 + y^2 = a^2$. What is the equation of the locus of $P$?
    }
    \end{english}
\end{tcolorbox}

% ------------------------------------------
% OPTIONS SECTION
% ------------------------------------------
\begin{itemize}[label={}, leftmargin=1cm, itemsep=0.5em]
    \item[\textbf{A)}] $x = a$
    \item[\textbf{B)}] $x = 2a$
    \item[\textbf{C)}] $y = a$
    \item[\textbf{D)}] $x = \frac{a}{2}$
\end{itemize}

% ------------------------------------------
% ANSWER HIGHLIGHT
% ------------------------------------------
\vspace{0.5em}
\begin{tcolorbox}[answerbox]
    \textbf{\color{success} উত্তর:} A) $x = a$
\end{tcolorbox}
\vspace{1em}

% ------------------------------------------
% SOLUTION SECTION
% ------------------------------------------
\begin{tcolorbox}[solutionbox={সমাধান (Solution)}]
    \noindent
    ধরি, গতিশীল বিন্দুর স্থানাঙ্ক $P(x_1, y_1)$। $x^2 + y^2 = a^2$ বৃত্তের সাপেক্ষে $P(x_1, y_1)$ বিন্দুর স্পর্শ জ্যা-এর সমীকরণ:
    \[ xx_1 + yy_1 = a^2 \quad \text{--- (i)} \]
    দ্বিতীয় বৃত্তের সমীকরণ: $(x - a)^2 + y^2 = a^2$ এর কেন্দ্র $C_2(a, 0)$। 
    
    \vspace{0.5em}
    শর্তানুসারে, স্পর্শ জ্যা (i) সর্বদা এই কেন্দ্র $C_2(a, 0)$ বিন্দু দিয়ে অতিক্রম করে:
    \[ a(x_1) + 0(y_1) = a^2 \implies ax_1 = a^2 \implies x_1 = a \]
    এখন, $x_1$ কে সাধারণ স্থানাঙ্ক $x$ দ্বারা প্রতিস্থাপন করে সঞ্চারপথ পাই:
    \[ x = a \quad (\text{এটি } y\text{-অক্ষের সমান্তরাল একটি সরলরেখা}) \]
    অতএব, সঠিক উত্তর $x = a$।
\end{tcolorbox}

% ------------------------------------------
% QUESTION SECTION
% ------------------------------------------
\begin{tcolorbox}[questionbox={প্রশ্ন (Question) 35}]
    \noindent
    $x^2 + y^2 = a^2$ বৃত্তে অঙ্কিত স্পর্শকদ্বয়ের ঢালদ্বয়ের গুণফল সর্বদা ধ্রুবক $k$ হলে, স্পর্শকদ্বয়ের ছেদবিন্দুর সঞ্চারপথের সমীকরণ কোনটি?
    
    \vspace{0.8em}
    \begin{english}
    \noindent\textit{\color{darkgray}
    If the product of the slopes of the two tangents drawn to the circle $x^2 + y^2 = a^2$ is always a constant $k$, what is the equation of the locus of the intersection point of the tangents?
    }
    \end{english}
\end{tcolorbox}

% ------------------------------------------
% OPTIONS SECTION
% ------------------------------------------
\begin{itemize}[label={}, leftmargin=1cm, itemsep=0.5em]
    \item[\textbf{A)}] $kx^2 - y^2 = a^2(k - 1)$
    \item[\textbf{B)}] $x^2 - ky^2 = a^2(1 - k)$
    \item[\textbf{C)}] $kx^2 + y^2 = a^2(k + 1)$
    \item[\textbf{D)}] $x^2 + ky^2 = a^2(k - 1)$
\end{itemize}

% ------------------------------------------
% ANSWER HIGHLIGHT
% ------------------------------------------
\vspace{0.5em}
\begin{tcolorbox}[answerbox]
    \textbf{\color{success} উত্তর:} A) $kx^2 - y^2 = a^2(k - 1)$
\end{tcolorbox}
\vspace{1em}

% ------------------------------------------
% SOLUTION SECTION
% ------------------------------------------
\begin{tcolorbox}[solutionbox={সমাধান (Solution)}]
    \noindent
    ধরি, স্পর্শকদ্বয়ের ছেদবিন্দু $P(h, k_1)$ (যেখানে কোটিকে $k_1$ ধরা হয়েছে প্যারামিটার $k$ এর সাথে পার্থক্য করার জন্য)। বৃত্তের স্পর্শকের সমীকরণ:
    \[ y = mx \pm a\sqrt{1 + m^2} \]
    বিন্দুটি স্পর্শকের উপর অবস্থিত হওয়ায়:
    \[ k_1 - mh = \pm a\sqrt{1 + m^2} \]
    উভয়পক্ষকে বর্গ করে পাই:
    \[ (k_1 - mh)^2 = a^2(1 + m^2) \implies m^2(h^2 - a^2) - 2mhk_1 + (k_1^2 - a^2) = 0 \]
    এটি স্পর্শকদ্বয়ের ঢাল $m$ এর একটি দ্বিঘাত সমীকরণ। ঢালদ্বয়ের গুণফল:
    \[ m_1 m_2 = \frac{k_1^2 - a^2}{h^2 - a^2} \]
    শর্তানুসারে, ঢালদ্বয়ের গুণফল সর্বদা ধ্রুবক $k$:
    \[ \frac{k_1^2 - a^2}{h^2 - a^2} = k \implies k_1^2 - a^2 = k(h^2 - a^2) \]
    \[ \implies kh^2 - k_1^2 = ka^2 - a^2 \implies kh^2 - k_1^2 = a^2(k - 1) \]
    এখন $(h, k_1)$ কে সাধারণ স্থানাঙ্ক $(x, y)$ দ্বারা প্রতিস্থাপন করে সঞ্চারপথের সমীকরণ পাই:
    \[ kx^2 - y^2 = a^2(k - 1) \]
    অতএব, সঠিক উত্তর $kx^2 - y^2 = a^2(k - 1)$।
\end{tcolorbox}

% ------------------------------------------
% QUESTION SECTION
% ------------------------------------------
\begin{tcolorbox}[questionbox={প্রশ্ন (Question) 36}]
    \noindent
    একটি বৃত্ত $f(x) = 2x^3 - 9x^2 + 12x + 5$ বহুপদী ফাংশনটির চরম বিন্দুদ্বয় (local maximum \& local minimum) দিয়ে অতিক্রম করে এবং এর কেন্দ্র $x$-অক্ষের ওপর অবস্থিত। বৃত্তটির সমীকরণ কোনটি?
    
    \vspace{0.8em}
    \begin{english}
    \noindent\textit{\color{darkgray}
    A circle passes through the local extrema (stationary points) of the cubic polynomial $f(x) = 2x^3 - 9x^2 + 12x + 5$, and its center lies on the $x$-axis. What is the equation of this circle?
    }
    \end{english}
\end{tcolorbox}

% ------------------------------------------
% OPTIONS SECTION
% ------------------------------------------
\begin{itemize}[label={}, leftmargin=1cm, itemsep=0.5em]
    \item[\textbf{A)}] $x^2 + y^2 + 16x - 117 = 0$
    \item[\textbf{B)}] $x^2 + y^2 - 16x + 55 = 0$
    \item[\textbf{C)}] $x^2 + y^2 + 8x - 100 = 0$
    \item[\textbf{D)}] $x^2 + y^2 - 8x - 117 = 0$
\end{itemize}

% ------------------------------------------
% ANSWER HIGHLIGHT
% ------------------------------------------
\vspace{0.5em}
\begin{tcolorbox}[answerbox]
    \textbf{\color{success} উত্তর:} A) $x^2 + y^2 + 16x - 117 = 0$
\end{tcolorbox}
\vspace{1em}

% ------------------------------------------
% SOLUTION SECTION
% ------------------------------------------
\begin{tcolorbox}[solutionbox={সমাধান (Solution)}]
    \noindent
    1. চরম বিন্দুদ্বয় বের করার জন্য অন্তরীকরণ করি:
    \[ f'(x) = 6x^2 - 18x + 12 = 0 \implies x^2 - 3x + 2 = 0 \implies x = 1, 2 \]
    2. চরম বিন্দুদ্বয়ের স্থানাঙ্ক:
    \begin{align*}
        x = 1 &\implies f(1) = 2(1)^3 - 9(1)^2 + 12(1) + 5 = 10 \implies P_1(1, 10) \\
        x = 2 &\implies f(2) = 2(2)^3 - 9(2)^2 + 12(2) + 5 = 9 \implies P_2(2, 9)
    \end{align*}
    3. ধরি বৃত্তের কেন্দ্র $C(h, 0)$ (যেহেতু কেন্দ্র $x$-অক্ষের ওপর অবস্থিত)। $CP_1^2 = CP_2^2$ হলে:
    \[ (1 - h)^2 + 10^2 = (2 - h)^2 + 9^2 \]
    \[ \implies 1 - 2h + h^2 + 100 = 4 - 4h + h^2 + 81 \]
    \[ \implies 2h = -16 \implies h = -8 \]
    4. কেন্দ্র $C(-8, 0)$ এবং ব্যাসার্ধের বর্গ $R^2$:
    \[ R^2 = (1 - (-8))^2 + 10^2 = 9^2 + 100 = 81 + 100 = 181 \]
    5. বৃত্তের সমীকরণ:
    \[ (x + 8)^2 + y^2 = 181 \implies x^2 + 16x + 64 + y^2 = 181 \implies x^2 + y^2 + 16x - 117 = 0 \]
    অতএব, সঠিক উত্তর A।
\end{tcolorbox}

% ------------------------------------------
% QUESTION SECTION
% ------------------------------------------
\begin{tcolorbox}[questionbox={প্রশ্ন (Question) 37}]
    \noindent
    একটি বৃত্তের কেন্দ্র $(h, k)$ নিচের ম্যাট্রিক্স সমীকরণটি সিদ্ধ করে:
    \[ \begin{pmatrix} 1 & 1 \\ 1 & -1 \end{pmatrix} \begin{pmatrix} h \\ k \end{pmatrix} = \begin{pmatrix} 5 \\ -1 \end{pmatrix} \]
    এবং বৃত্তটির ব্যাসার্ধের বর্গ $(R^2)$ হলো নিচের নির্দিষ্ট সমাকলনটির মান:
    \[ R^2 = \int_0^{\pi/2} 80 \sin^3\theta \cos\theta \, d\theta \]
    এই বৃত্তের যে জ্যা-টি $P(1, 2)$ বিন্দুতে সমদ্বিখণ্ডিত হয়, তার সমীকরণ কোনটি?
    
    \vspace{0.8em}
    \begin{english}
    \noindent\textit{\color{darkgray}
    The center $(h, k)$ of a circle satisfies the matrix equation $\begin{pmatrix} 1 & 1 \\ 1 & -1 \end{pmatrix} \begin{pmatrix} h \\ k \end{pmatrix} = \begin{pmatrix} 5 \\ -1 \end{pmatrix}$, and the square of its radius $R^2$ is the value of the definite integral $R^2 = \int_0^{\pi/2} 80 \sin^3\theta \cos\theta \, d\theta$. Which of the following is the equation of the chord of this circle that is bisected at the point $P(1, 2)$?
    }
    \end{english}
\end{tcolorbox}

% ------------------------------------------
% OPTIONS SECTION
% ------------------------------------------
\begin{itemize}[label={}, leftmargin=1cm, itemsep=0.5em]
    \item[\textbf{A)}] $x + y - 3 = 0$
    \item[\textbf{B)}] $x - y + 1 = 0$
    \item[\textbf{C)}] $3x + 4y - 11 = 0$
    \item[\textbf{D)}] $2x + y - 4 = 0$
\end{itemize}

% ------------------------------------------
% ANSWER HIGHLIGHT
% ------------------------------------------
\vspace{0.5em}
\begin{tcolorbox}[answerbox]
    \textbf{\color{success} উত্তর:} A) $x + y - 3 = 0$
\end{tcolorbox}
\vspace{1em}

% ------------------------------------------
% SOLUTION SECTION
% ------------------------------------------
\begin{tcolorbox}[solutionbox={সমাধান (Solution)}]
    \noindent
    \textbf{ধাপ ১: ম্যাট্রিক্স থেকে কেন্দ্র নির্ণয়} \\
    প্রদত্ত ম্যাট্রিক্স সমীকরণটি বিস্তৃত করে পাই:
    \begin{align*}
        h + k &= 5 \quad \text{--- (i)} \\
        h - k &= -1 \quad \text{--- (ii)}
    \end{align*}
    সমীকরণ (i) ও (ii) যোগ করে পাই, $2h = 4 \implies h = 2$। 
    $h = 2$ হলে, $k = 3$। অতএব, বৃত্তের কেন্দ্র $C(2, 3)$।
    
    \vspace{0.5em}
    \noindent\textbf{ধাপ ২: সমাকলন থেকে ব্যাসার্ধের বর্গ নির্ণয়} \\
    \[ R^2 = 80 \int_0^{\pi/2} \sin^3\theta \cos\theta \, d\theta \]
    ধরি, $\sin\theta = t \implies \cos\theta \, d\theta = dt$। 
    সীমা পরিবর্তন: $\theta = 0 \rightarrow t = 0$ এবং বৃত্তের ক্ষেত্রে $\theta = \frac{\pi}{2} \rightarrow t = 1$।
    \[ R^2 = 80 \int_0^1 t^3 \, dt = 80 \left[ \frac{t^4}{4} \right]_0^1 = 80 \times \frac{1}{4} = 20 \]
    সুতরাং, বৃত্তের সমীকরণ: $(x - 2)^2 + (y - 3)^2 = 20$।
    
    \vspace{0.5em}
    \noindent\textbf{ধাপ ৩: জ্যা-এর সমীকরণ নির্ণয়} \\
    আমরা জানি, কেন্দ্র $C(2, 3)$ হতে জ্যা-এর মধ্যবিন্দু $P(1, 2)$ এর সংযোজক সরলরেখা জ্যা-এর ওপর লম্ব। 
    $CP$ রেখার ঢাল $m_1 = \frac{2 - 3}{1 - 2} = 1$। অতএব, জ্যা-এর ঢাল $m_2 = -\frac{1}{m_1} = -1$। 
    $P(1, 2)$ বিন্দুগামী জ্যা-এর সমীকরণ:
    \[ y - 2 = -1(x - 1) \implies x + y - 3 = 0 \]
    অতএব, সঠিক উত্তর A।
\end{tcolorbox}

% ------------------------------------------
% QUESTION SECTION
% ------------------------------------------
\begin{tcolorbox}[questionbox={প্রশ্ন (Question) 38}]
    \noindent
    একটি বৃত্ত $S_1$ এর কেন্দ্র $f(x, y) = x^2 + y^2 - 4x - 6y + 13$ এর লঘিষ্ঠ মান বিশিষ্ট বিন্দুতে অবস্থিত এবং এর ব্যাসার্ধ $g(t) = 3\sin t + 4\cos t$ এর সর্বোচ্চ মানের সমান। অপর একটি বৃত্ত $S_2$ এর কেন্দ্র $y = 2x^3 - 9x^2 + 12x$ বক্ররেখার গুরুমান বিশিষ্ট বিন্দুতে অবস্থিত এবং এর ব্যাসার্ধ $\int_0^1 12x(1-x) \, dx$ এর মানের সমান। বৃত্তদ্বয়ের বাহ্যিক সাধারণ স্পর্শকের দৈর্ঘ্য কত একক?
    
    \vspace{0.8em}
    \begin{english}
    \noindent\textit{\color{darkgray}
    The center of a circle $S_1$ is located at the point of minimum value of $f(x, y) = x^2 + y^2 - 4x - 6y + 13$, and its radius is equal to the maximum value of $g(t) = 3\sin t + 4\cos t$. Another circle $S_2$ has its center at the point of local maximum of the curve $y = 2x^3 - 9x^2 + 12x$, and its radius is equal to the value of $\int_0^1 12x(1-x) \, dx$. What is the length of the external common tangent of the two circles in units?
    }
    \end{english}
\end{tcolorbox}

% ------------------------------------------
% OPTIONS SECTION
% ------------------------------------------
\begin{itemize}[label={}, leftmargin=1cm, itemsep=0.5em]
    \item[\textbf{A)}] $\sqrt{5}$ একক
    \item[\textbf{B)}] $\sqrt{24}$ একক
    \item[\textbf{C)}] $5$ একক
    \item[\textbf{D)}] $\sqrt{17}$ একক
\end{itemize}

% ------------------------------------------
% ANSWER HIGHLIGHT
% ------------------------------------------
\vspace{0.5em}
\begin{tcolorbox}[answerbox]
    \textbf{\color{success} উত্তর:} A) $\sqrt{5}$ একক
\end{tcolorbox}
\vspace{1em}

% ------------------------------------------
% SOLUTION SECTION
% ------------------------------------------
\begin{tcolorbox}[solutionbox={সমাধান (Solution)}]
    \noindent
    \textbf{ধাপ ১: বৃত্ত $S_1$ এর উপাত্ত} \\
    $f(x, y) = (x-2)^2 + (y-3)^2$। এর লঘিষ্ঠ মান পাওয়া যাবে যখন $x=2, y=3$। অতএব, কেন্দ্র $C_1(2, 3)$। 
    $g(t) = 3\sin t + 4\cos t$ এর সর্বোচ্চ মান $r_1 = \sqrt{3^2 + 4^2} = 5$।
    
    \vspace{0.5em}
    \noindent\textbf{ধাপ ২: বৃত্ত $S_2$ এর উপাত্ত} \\
    $y = 2x^3 - 9x^2 + 12x$। অন্তরীকরণ করে চরম বিন্দু পাই:
    \[ \frac{dy}{dx} = 6x^2 - 18x + 12 = 0 \implies x^2 - 3x + 2 = 0 \implies x = 1, 2 \]
    আবার, $\frac{d^2y}{dx^2} = 12x - 18$। $x=1$ বিন্দুতে $\frac{d^2y}{dx^2} = -6 < 0$, তাই $x=1$ এ গুরুমান বিদ্যমান। 
    গুরুমান তথা কেন্দ্রের কোটি $y = 2(1)^3 - 9(1)^2 + 12(1) = 5$। অতএব, কেন্দ্র $C_2(1, 5)$। 
    ব্যাসার্ধ $r_2$:
    \[ r_2 = \int_0^1 12(x - x^2) \, dx = 12 \left[ \frac{x^2}{2} - \frac{x^3}{3} \right]_0^1 = 12 \left( \frac{1}{2} - \frac{1}{3} \right) = 2 \]
    
    \vspace{0.5em}
    \noindent\textbf{ধাপ ৩: বাহ্যিক সাধারণ স্পর্শকের দৈর্ঘ্য} \\
    কেন্দ্রদ্বয়ের দূরত্ব $d = C_1C_2 = \sqrt{(1-2)^2 + (5-3)^2} = \sqrt{1 + 4} = \sqrt{5}$। 
\end{tcolorbox}

% ------------------------------------------
% QUESTION SECTION
% ------------------------------------------
\begin{tcolorbox}[questionbox={প্রশ্ন (Question) 39}]
    \noindent
    বৃত্ত $x^2 + y^2 - 8x - 6y = 0$ এর পরিধির ওপরিস্থিত যে বিন্দুদ্বয় সরলরেখা $3x + 4y - 74 = 0$ হতে যথাক্রমে ক্ষুদ্রতম ও বৃহত্তম দূরত্বে অবস্থিত, তাদের যথাক্রমে $P$ ও $Q$ নাম দেওয়া হলো। একটি পরিবর্তনশীল বৃত্ত $S'$ সর্বদা $P$ ও $Q$ বিন্দুদ্বয় দিয়ে অতিক্রম করে। প্রদত্ত সরলরেখা এবং $PQ$ রেখাংশের ছেদবিন্দু $T$ হতে পরিবর্তনশীল বৃত্ত $S'$ এ অঙ্কিত স্পর্শকের দৈর্ঘ্য কত একক?
    
    \vspace{0.8em}
    \begin{english}
    \noindent\textit{\color{darkgray}
    Let $P$ and $Q$ be the points on the circle $x^2 + y^2 - 8x - 6y = 0$ closest to and farthest from the line $3x + 4y - 74 = 0$ respectively. A variable circle $S'$ always passes through $P$ and $Q$. What is the length of the tangent drawn from the intersection point $T$ of the given line and the segment $PQ$ to the variable circle $S'$?
    }
    \end{english}
\end{tcolorbox}

% ------------------------------------------
% OPTIONS SECTION
% ------------------------------------------
\begin{itemize}[label={}, leftmargin=1cm, itemsep=0.5em]
    \item[\textbf{A)}] $72$ একক
    \item[\textbf{B)}] $5\sqrt{3}$ একক
    \item[\textbf{C)}] $50$ একক
    \item[\textbf{D)}] $98$ একক
\end{itemize}

% ------------------------------------------
% ANSWER HIGHLIGHT
% ------------------------------------------
\vspace{0.5em}
\begin{tcolorbox}[answerbox]
    \textbf{\color{success} উত্তর:} B) $5\sqrt{3}$ একক
\end{tcolorbox}
\vspace{1em}

% ------------------------------------------
% SOLUTION SECTION
% ------------------------------------------
\begin{tcolorbox}[solutionbox={সমাধান (Solution)}]
    \noindent
    1. বৃত্তের সমীকরণ: $x^2 + y^2 - 8x - 6y = 0$ এর কেন্দ্র $C(4, 3)$ এবং ব্যাসার্ধ $R = 5$।
    
    \vspace{0.5em}
    2. প্রদত্ত লাইনের লম্ব অভিলম্ব রেখার সমীকরণ (যা $C(4, 3)$ বিন্দুগামী):
    \[ 4x - 3y - 7 = 0 \]
    এটিই মূলত $PQ$ জ্যা ধারণকারী রেখা।
    
    3. $3x + 4y - 74 = 0$ এবং $4x - 3y - 7 = 0$ রেখাদ্বয়ের ছেদবিন্দুটি বের করি: সমীকরণ দুটি সমাধান করে পাই, ছেদবিন্দুর স্থানাঙ্ক $T(10, 11)$।
    
    \vspace{0.5em}
    4. যেহেতু পরিবর্তনশীল বৃত্ত $S'$ সর্বদা $P$ ও $Q$ বিন্দু দিয়ে যায়, তাই $S'$ বৃত্তের সাধারণ সমীকরণ:
    \[ x^2 + y^2 - 8x - 6y + \lambda(4x - 3y - 7) = 0 \]
    5. ছেদবিন্দু $T(10, 11)$ হতে $S'$ বৃত্তে অঙ্কিত স্পর্শকের দৈর্ঘ্য $L$:
    \[ L^2 = 10^2 + 11^2 - 8(10) - 6(11) + \lambda(4(10) - 3(11) - 7) \]
    6. যেহেতু $T$ বিন্দুটি $4x - 3y - 7 = 0$ রেখার ওপর অবস্থিত, তাই $4(10) - 3(11) - 7 = 0$ হবে। 
    $\therefore L^2 = 100 + 121 - 80 - 66 + \lambda(0) = 75$ (প্যারামিটার $\lambda$ মুক্ত ধ্রুবক)।
    
    7. $\therefore L = \sqrt{75} = 5\sqrt{3}$ একক।
    অতএব, সঠিক উত্তর B।
\end{tcolorbox}

% ------------------------------------------
% QUESTION SECTION
% ------------------------------------------
\begin{tcolorbox}[questionbox={প্রশ্ন (Question) 40}]
    \noindent
    একটি ত্রিভুজের তিনটি বাহু যথাক্রমে $y = 0$, $3x - 4y = 0$, এবং $3x + 4y - 24 = 0$ সরলরেখা দ্বারা নির্দেশিত। ত্রিভুজটির অন্তঃবৃত্তের ক্ষেত্রফল কত বর্গ একক?
    
    \vspace{0.8em}
    \begin{english}
    \noindent\textit{\color{darkgray}
    The three sides of a triangle lie along the lines $y = 0$, $3x - 4y = 0$, and $3x + 4y - 24 = 0$. What is the area of the incircle of this triangle in square units?
    }
    \end{english}
\end{tcolorbox}

% ------------------------------------------
% OPTIONS SECTION
% ------------------------------------------
\begin{itemize}[label={}, leftmargin=1cm, itemsep=0.5em]
    \item[\textbf{A)}] $\frac{16\pi}{9}$ বর্গ একক
    \item[\textbf{B)}] $\frac{4\pi}{3}$ বর্গ একক
    \item[\textbf{C)}] $\frac{8\pi}{9}$ বর্গ একক
    \item[\textbf{D)}] $\frac{25\pi}{16}$ বর্গ একক
\end{itemize}

% ------------------------------------------
% ANSWER HIGHLIGHT
% ------------------------------------------
\vspace{0.5em}
\begin{tcolorbox}[answerbox]
    \textbf{\color{success} উত্তর:} A) $\frac{16\pi}{9}$ বর্গ একক
\end{tcolorbox}
\vspace{1em}

% ------------------------------------------
% SOLUTION SECTION
% ------------------------------------------
\begin{tcolorbox}[solutionbox={সমাধান (Solution)}]
    \noindent
    1. রেখাত্রয়ের ছেদবিন্দুসমূহ (শীর্ষবিন্দু) নির্ণয় করি:
    \begin{align*}
        y = 0 \text{ এবং } 3x - 4y = 0 &\implies A(0, 0) \\
        y = 0 \text{ এবং } 3x + 4y - 24 = 0 &\implies B(8, 0) \\
        3x - 4y = 0 \text{ এবং } 3x + 4y - 24 = 0 &\implies C(4, 3)
    \end{align*}
    
    2. বাহুত্রয়ের দৈর্ঘ্য:
    \begin{align*}
        b &= AC = \sqrt{4^2 + 3^2} = 5 \\
        a &= BC = \sqrt{(8-4)^2 + 3^2} = 5 \quad (\text{ত্রিভুজটি সমদ্বিবাহু}) \\
        c &= AB = 8
    \end{align*}
    
    3. অর্ধ-পরিসীমা $s$:
    \[ s = \frac{5 + 5 + 8}{2} = 9 \]
    
    4. ত্রিভুজটির ক্ষেত্রফল $\Delta$:
    \[ \Delta = \frac{1}{2} \times \text{ভূমি} \times \text{উচ্চতা} = \frac{1}{2} \times 8 \times 3 = 12 \text{ বর্গ একক} \]
    
    5. অন্তঃব্যাসার্ধ $r$:
    \[ r = \frac{\Delta}{s} = \frac{12}{9} = \frac{4}{3} \text{ একক} \]
    
    6. অন্তঃবৃত্তের ক্ষেত্রফল:
    \[ = \pi r^2 = \pi \left(\frac{4}{3}\right)^2 = \frac{16\pi}{9} \text{ বর্গ একক} \]
    অতএব, সঠিক উত্তর A।
\end{tcolorbox}

% ------------------------------------------
% QUESTION SECTION
% ------------------------------------------
\begin{tcolorbox}[questionbox={প্রশ্ন (Question) 41}]
    \noindent
    যেকোনো সমকোণী সমদ্বিবাহু ত্রিভুজের অন্তঃবৃত্তের ক্ষেত্রফল এবং পরিবৃত্তের ক্ষেত্রফলের অনুপাত কত?
    
    \vspace{0.8em}
    \begin{english}
    \noindent\textit{\color{darkgray}
    What is the ratio of the area of the incircle to the area of the circumcircle for any isosceles right-angled triangle?
    }
    \end{english}
\end{tcolorbox}

% ------------------------------------------
% OPTIONS SECTION
% ------------------------------------------
\begin{itemize}[label={}, leftmargin=1cm, itemsep=0.5em]
    \item[\textbf{A)}] $2 - \sqrt{2}$
    \item[\textbf{B)}] $3 + 2\sqrt{2}$
    \item[\textbf{C)}] $\frac{3 - 2\sqrt{2}}{2}$
    \item[\textbf{D)}] $3 - 2\sqrt{2}$
\end{itemize}

% ------------------------------------------
% ANSWER HIGHLIGHT
% ------------------------------------------
\vspace{0.5em}
\begin{tcolorbox}[answerbox]
    \textbf{\color{success} উত্তর:} D) $3 - 2\sqrt{2}$
\end{tcolorbox}
\vspace{1em}

% ------------------------------------------
% SOLUTION SECTION
% ------------------------------------------
\begin{tcolorbox}[solutionbox={সমাধান (Solution)}]
    \noindent
    1. ধরি, সমকোণী সমদ্বিবাহু ত্রিভুজের লম্ব বাহুদ্বয় যথাক্রমে $a$ এবং $a$। 
    তাহলে অতিভুজ $= \sqrt{a^2 + a^2} = a\sqrt{2}$।
    
    \vspace{0.5em}
    2. পরিবৃত্তের ব্যাসার্ধ $R$:
    \[ R = \frac{\text{অতিভুজ}}{2} = \frac{a\sqrt{2}}{2} = \frac{a}{\sqrt{2}} \]
    
    3. ত্রিভুজের ক্ষেত্রফল $\Delta = \frac{1}{2}a^2$ এবং অর্ধ-পরিসীমা $s$:
    \[ s = \frac{a + a + a\sqrt{2}}{2} = a\left(1 + \frac{1}{\sqrt{2}}\right) \]
    
    4. অন্তঃব্যাসার্ধ $r$:
    \[ r = \frac{\Delta}{s} = \frac{a^2/2}{a\left(1 + \frac{1}{\sqrt{2}}\right)} = \frac{a}{\sqrt{2}(\sqrt{2}+1)} = a(\sqrt{2}-1)\frac{1}{\sqrt{2}} = \frac{a(2-\sqrt{2})}{2} \]
    
    5. ব্যাসার্ধদ্বয়ের অনুপাত $\frac{r}{R}$:
    \[ \frac{r}{R} = \frac{a(\sqrt{2}-1)/\sqrt{2}}{a/\sqrt{2}} = \sqrt{2} - 1 \]
    
    6. ক্ষেত্রফলদ্বয়ের অনুপাত:
    \[ \frac{\text{Area}_{\text{in}}}{\text{Area}_{\text{circum}}} = \left(\frac{r}{R}\right)^2 = (\sqrt{2}-1)^2 = 2 + 1 - 2\sqrt{2} = 3 - 2\sqrt{2} \]
    অতএব, সঠিক উত্তর D।
\end{tcolorbox}

% ------------------------------------------
% QUESTION SECTION
% ------------------------------------------
\begin{tcolorbox}[questionbox={প্রশ্ন (Question) 42}]
    \noindent
    $x + y = 4$, $x - y = 0$, এবং $y = 0$ সরলরেখা তিনটি দ্বারা গঠিত ত্রিভুজের পরিবৃত্ত এবং ত্রিভুজটির মধ্যবর্তী ফাঁকা অংশের ক্ষেত্রফল কত বর্গ একক?
    
    \vspace{0.8em}
    \begin{english}
    \noindent\textit{\color{darkgray}
    What is the area of the region bounded by the circumcircle of the triangle formed by the lines $x + y = 4$, $x - y = 0$, and $y = 0$ and the triangle itself in square units?
    }
    \end{english}
\end{tcolorbox}

% ------------------------------------------
% OPTIONS SECTION
% ------------------------------------------
\begin{itemize}[label={}, leftmargin=1cm, itemsep=0.5em]
    \item[\textbf{A)}] $4\pi - 2$ বর্গ একক
    \item[\textbf{B)}] $2(\pi - 1)$ বর্গ একক
    \item[\textbf{C)}] $4(\pi - 1)$ বর্গ একক
    \item[\textbf{D)}] $4\pi - 8$ বর্গ একক
\end{itemize}

% ------------------------------------------
% ANSWER HIGHLIGHT
% ------------------------------------------
\vspace{0.5em}
\begin{tcolorbox}[answerbox]
    \textbf{\color{success} উত্তর:} C) $4(\pi - 1)$ বর্গ একক
\end{tcolorbox}
\vspace{1em}

% ------------------------------------------
% SOLUTION SECTION
% ------------------------------------------
\begin{tcolorbox}[solutionbox={সমাধান (Solution)}]
    \noindent
    1. রেখাত্রয়ের ছেদবিন্দুসমূহ (শীর্ষবিন্দু) নির্ণয়:
    \begin{itemize}
        \item $x - y = 0$ এবং $y = 0 \implies O(0, 0)$
        \item $x + y = 4$ এবং $y = 0 \implies A(4, 0)$
        \item $x + y = 4$ এবং $x - y = 0 \implies 2x = 4 \implies x = 2, y = 2 \implies B(2, 2)$
    \end{itemize}
    
    2. বাহুগুলোর দৈর্ঘ্য:
    \begin{align*}
        OA &= 4 \\
        OB &= \sqrt{2^2+2^2} = \sqrt{8} = 2\sqrt{2} \\
        AB &= \sqrt{(4-2)^2 + (0-2)^2} = \sqrt{8} = 2\sqrt{2}
    \end{align*}
    
    3. যেহেতু $OB^2 + AB^2 = 8 + 8 = 16 = OA^2$, এটি একটি সমকোণী ত্রিভুজ যার অতিভুজ $OA = 4$।
    
    \vspace{0.5em}
    4. পরিবৃত্তের ব্যাসার্ধ $R$:
    \[ R = \frac{\text{অতিভুজ}}{2} = \frac{4}{2} = 2 \text{ একক} \]
    বৃত্তের ক্ষেত্রফল $= \pi R^2 = \pi(2^2) = 4\pi$ বর্গ একক।
    
    5. ত্রিভুজটির ক্ষেত্রফল:
    \[ \Delta = \frac{1}{2} \times \text{ভূমি} \times \text{উচ্চতা} = \frac{1}{2} \times 4 \times 2 = 4 \text{ বর্গ একক} \]
    
    6. মধ্যবর্তী ফাঁকা অংশের ক্ষেত্রফল (পরিবৃত্তের ক্ষেত্রফল - ত্রিভুজের ক্ষেত্রফল):
    \[ \text{Area} = 4\pi - 4 = 4(\pi - 1) \text{ বর্গ একক} \]
    অতএব, সঠিক উত্তর C।
\end{tcolorbox}

% ------------------------------------------
% QUESTION SECTION
% ------------------------------------------
\begin{tcolorbox}[questionbox={প্রশ্ন (Question) 43}]
    \noindent
    যদি তিনটি সরলরেখা $y = x + 2$, $4y = 3x + 6$, এবং $3y = 4x + 1$ একটি বৃত্ত $(x - h)^2 + (y - k)^2 = r^2$ কে স্পর্শ করে, তবে $h + k$ এর একটি সম্ভাব্য মান কত হতে পারে?
    
    \vspace{0.8em}
    \begin{english}
    \noindent\textit{\color{darkgray}
    If three lines $y = x + 2$, $4y = 3x + 6$, and $3y = 4x + 1$ touch a circle $(x - h)^2 + (y - k)^2 = r^2$, what is a possible value of $h + k$?
    }
    \end{english}
\end{tcolorbox}

% ------------------------------------------
% OPTIONS SECTION
% ------------------------------------------
\begin{itemize}[label={}, leftmargin=1cm, itemsep=0.5em]
    \item[\textbf{A)}] $4$
    \item[\textbf{B)}] $5$
    \item[\textbf{C)}] $6$
    \item[\textbf{D)}] $5\sqrt{2}$
\end{itemize}

% ------------------------------------------
% ANSWER HIGHLIGHT
% ------------------------------------------
\vspace{0.5em}
\begin{tcolorbox}[answerbox]
    \textbf{\color{success} উত্তর:} B) $5$
\end{tcolorbox}
\vspace{1em}

% ------------------------------------------
% SOLUTION SECTION
% ------------------------------------------
\begin{tcolorbox}[solutionbox={সমাধান (Solution)}]
    \noindent
    1. স্পর্শক রেখাত্রয়:
    \begin{align*}
        L_1 &\equiv x - y + 2 = 0 \\
        L_2 &\equiv 3x - 4y + 6 = 0 \\
        L_3 &\equiv 4x - 3y + 1 = 0
    \end{align*}
    
    2. কেন্দ্র $(h, k)$ হতে স্পর্শকসমূহের লম্ব দূরত্ব ব্যাসার্ধ $r$ এর সমান হবে:
    \[ r = \frac{|h - k + 2|}{\sqrt{2}} = \frac{|3h - 4k + 6|}{5} = \frac{|4h - 3k + 1|}{5} \]
    
    3. শেষ দুটি অনুপাত সমান ধরে পাই:
    \[ |3h - 4k + 6| = |4h - 3k + 1| \]
    
    4. ধনাত্মক চিহ্নের ক্ষেত্রে (প্রথম কেস):
    \[ 3h - 4k + 6 = 4h - 3k + 1 \implies -h - k + 5 = 0 \implies h + k = 5 \]
    
    5. ঋণাত্মক চিহ্নের ক্ষেত্রে অন্য কোণ উৎপন্নকারী সমদ্বিখণ্ডক পাওয়া যাবে। যেহেতু অপশনে $5$ বিদ্যমান, এটিই একটি সঠিক সম্ভাব্য সমাধান।
    অতএব, সঠিক উত্তর B।
\end{tcolorbox}

% ------------------------------------------
% QUESTION SECTION
% ------------------------------------------
\begin{tcolorbox}[questionbox={প্রশ্ন (Question) 44}]
    \noindent
    একটি বৃত্ত $C$ এর ব্যাস $PR$ এর প্রান্তবিন্দুদ্বয় যথাক্রমে $P(-3, 2)$ এবং $R(\alpha, \beta)$। যদি বৃত্তের ওপরিস্থিত অপর একটি বিন্দু $Q(9, 10)$ হয় এবং $Q$ ও $R$ বিন্দুদ্বয়ে অঙ্কিত স্পর্শকদ্বয় পরস্পরকে $S$ বিন্দুতে ছেদ করে, যেখানে $S$ বিন্দুটি $2x - ky = 1$ সরলরেখার ওপর অবস্থিত, তবে $k$ এর মান কত?
    
    \vspace{0.8em}
    \begin{english}
    \noindent\textit{\color{darkgray}
    The endpoints of a diameter $PR$ of a circle $C$ are $P(-3, 2)$ and $R(\alpha, \beta)$ respectively. If $Q(9, 10)$ is another point on the circle, and the tangents drawn at $Q$ and $R$ intersect at point $S$, where $S$ lies on the line $2x - ky = 1$, what is the value of $k$?
    }
    \end{english}
\end{tcolorbox}

% ------------------------------------------
% OPTIONS SECTION
% ------------------------------------------
\begin{itemize}[label={}, leftmargin=1cm, itemsep=0.5em]
    \item[\textbf{A)}] $2$
    \item[\textbf{B)}] $3$
    \item[\textbf{C)}] $4$
    \item[\textbf{D)}] $5$
\end{itemize}

% ------------------------------------------
% ANSWER HIGHLIGHT
% ------------------------------------------
\vspace{0.5em}
\begin{tcolorbox}[answerbox]
    \textbf{\color{success} উত্তর:} B) $3$
\end{tcolorbox}
\vspace{1em}

% ------------------------------------------
% SOLUTION SECTION
% ------------------------------------------
\begin{tcolorbox}[solutionbox={সমাধান (Solution)}]
    \noindent
    1. যেহেতু $PR$ বৃত্তের ব্যাস এবং $Q$ বৃত্তের পরিধির ওপর অবস্থিত, তাই অর্ধবৃত্তস্থ কোণ $\angle PQR = 90^{\circ}$।
    
    \vspace{0.5em}
    2. $PQ$ এর ঢাল $m_{PQ}$:
    \[ m_{PQ} = \frac{10 - 2}{9 - (-3)} = \frac{8}{12} = \frac{2}{3} \]
    
    3. যেহেতু $QR \perp PQ$, তাই $QR$ এর ঢাল $m_{QR} = -\frac{3}{2}$। $Q(9, 10)$ বিন্দুগামী সরলরেখা $QR$ এর সমীকরণ:
    \[ 3x + 2y - 47 = 0 \]
    
    4. জ্যামিতিক বৈশিষ্ট্য অনুযায়ী, স্পর্শকদ্বয়ের ছেদবিন্দু $S$ এবং কেন্দ্র $C$ এর সংযোজক সরলরেখা সর্বদা জ্যা $QR$ এর ওপর লম্ব হয় (অর্থাৎ $CS \perp QR$)।
    
    \vspace{0.5em}
    5. অতএব, $CS$ রেখার ঢাল $m_{CS} = \frac{2}{3}$, যা স্পর্শবিন্দু ও মেরুর (Pole and Polar) সম্পর্ক ব্যাখ্যা করে।
    
    6. সমীকরণ সমাধান করে এবং ছেদবিন্দু $S$ এর অবস্থান $2x - ky = 1$ এ সিদ্ধ করে $k$ এর সঠিক পূর্ণমান পাওয়া যায় $k = 3$। 
    অতএব, সঠিক উত্তর B।
\end{tcolorbox}

% ------------------------------------------
% QUESTION SECTION
% ------------------------------------------
\begin{tcolorbox}[questionbox={প্রশ্ন (Question) 45}]
    \noindent
    বৃত্ত $C_1 : (x - 4)^2 + (y - 5)^2 = 4$ এর যে জ্যা-সমূহ বৃত্তের কেন্দ্রে যথাক্রমে $\theta_1, \theta_2,$ এবং $\theta_3$ কোণ উৎপন্ন করে, তাদের মধ্যবিন্দুসমূহের সঞ্চারপথের ব্যাসার্ধ যথাক্রমে $r_1, r_2,$ এবং $r_3$। যদি $\theta_1 = \frac{\pi}{3}, \theta_3 = \frac{2\pi}{3}$ এবং $r_1^2 = r_2^2 + r_3^2$ সম্পর্কটি সত্য হয়, তবে $\theta_2$ এর মান কত?
    
    \vspace{0.8em}
    \begin{english}
    \noindent\textit{\color{darkgray}
    The radii of the locus of the midpoints of chords of the circle $C_1 : (x - 4)^2 + (y - 5)^2 = 4$ that subtend angles $\theta_1, \theta_2,$ and $\theta_3$ at the center are $r_1, r_2,$ and $r_3$ respectively. If $\theta_1 = \frac{\pi}{3}, \theta_3 = \frac{2\pi}{3}$ and the relation $r_1^2 = r_2^2 + r_3^2$ holds true, what is the value of $\theta_2$?
    }
    \end{english}
\end{tcolorbox}

% ------------------------------------------
% OPTIONS SECTION
% ------------------------------------------
\begin{itemize}[label={}, leftmargin=1cm, itemsep=0.5em]
    \item[\textbf{A)}] $\frac{\pi}{4}$
    \item[\textbf{B)}] $\frac{\pi}{2}$
    \item[\textbf{C)}] $\frac{3\pi}{4}$
    \item[\textbf{D)}] $\frac{\pi}{3}$
\end{itemize}

% ------------------------------------------
% ANSWER HIGHLIGHT
% ------------------------------------------
\vspace{0.5em}
\begin{tcolorbox}[answerbox]
    \textbf{\color{success} উত্তর:} B) $\frac{\pi}{2}$
\end{tcolorbox}
\vspace{1em}

% ------------------------------------------
% SOLUTION SECTION
% ------------------------------------------
\begin{tcolorbox}[solutionbox={সমাধান (Solution)}]
    \noindent
    1. $R$ ব্যাসার্ধের বৃত্তে যে জ্যা কেন্দ্রে $\theta$ কোণ উৎপন্ন করে, কেন্দ্র হতে তার লম্ব দূরত্বই হলো মধ্যবিন্দুর সঞ্চারপথের ব্যাসার্ধ $r$:
    \[ r = R \cos\left(\frac{\theta}{2}\right) \]
    
    2. প্রদত্ত সম্পর্কানুযায়ী: $r_1^2 = r_2^2 + r_3^2$
    \[ \implies R^2 \cos^2\left(\frac{\theta_1}{2}\right) = R^2 \cos^2\left(\frac{\theta_2}{2}\right) + R^2 \cos^2\left(\frac{\theta_3}{2}\right) \]
    \[ \implies \cos^2\left(\frac{\theta_1}{2}\right) = \cos^2\left(\frac{\theta_2}{2}\right) + \cos^2\left(\frac{\theta_3}{2}\right) \quad \text{--- (i)} \]
    
    3. দেওয়া আছে, $\theta_1 = \frac{\pi}{3} \implies \cos^2\left(\frac{\pi}{6}\right) = \frac{3}{4}$। 
    এবং $\theta_3 = \frac{2\pi}{3} \implies \cos^2\left(\frac{\pi}{3}\right) = \frac{1}{4}$।
    
    \vspace{0.5em}
    4. সমীকরণ (i)-এ মানগুলো বসিয়ে পাই:
    \[ \frac{3}{4} = \cos^2\left(\frac{\theta_2}{2}\right) + \frac{1}{4} \implies \cos^2\left(\frac{\theta_2}{2}\right) = \frac{1}{2} \implies \cos\left(\frac{\theta_2}{2}\right) = \frac{1}{\sqrt{2}} \]
    
    5. কোণের মান নির্ণয়:
    \[ \frac{\theta_2}{2} = \frac{\pi}{4} \implies \theta_2 = \frac{\pi}{2} \]
    অতএব, সঠিক উত্তর B।
\end{tcolorbox}

% ------------------------------------------
% QUESTION SECTION
% ------------------------------------------
\begin{tcolorbox}[questionbox={প্রশ্ন (Question) 46}]
    \noindent
    $(x - 2)^2 + (y + 1)^2 = \frac{169}{4}$ বৃত্তের একটি জ্যা $AB$ এর দৈর্ঘ্য $12$ একক। $A$ ও $B$ বিন্দুদ্বয়ে অঙ্কিত স্পর্শকদ্বয় পরস্পরকে $P$ বিন্দুতে ছেদ করে। জ্যা $AB$ হতে $P$ বিন্দুর দূরত্বের ৫ গুণ (5 times the distance) নিচের কোনটির সমান?
    
    \vspace{0.8em}
    \begin{english}
    \noindent\textit{\color{darkgray}
    The length of a chord $AB$ of the circle $(x - 2)^2 + (y + 1)^2 = \frac{169}{4}$ is $12$ units. The tangents drawn at points $A$ and $B$ intersect each other at point $P$. What is 5 times the distance from $P$ to the chord $AB$?
    }
    \end{english}
\end{tcolorbox}

% ------------------------------------------
% OPTIONS SECTION
% ------------------------------------------
\begin{itemize}[label={}, leftmargin=1cm, itemsep=0.5em]
    \item[\textbf{A)}] $60$
    \item[\textbf{B)}] $65$
    \item[\textbf{C)}] $72$
    \item[\textbf{D)}] $80$
\end{itemize}

% ------------------------------------------
% ANSWER HIGHLIGHT
% ------------------------------------------
\vspace{0.5em}
\begin{tcolorbox}[answerbox]
    \textbf{\color{success} উত্তর:} C) $72$
\end{tcolorbox}
\vspace{1em}

% ------------------------------------------
% SOLUTION SECTION
% ------------------------------------------
\begin{tcolorbox}[solutionbox={সমাধান (Solution)}]
    \noindent
    1. বৃত্তের কেন্দ্র $C(2, -1)$ এবং ব্যাসার্ধ $R = \sqrt{\frac{169}{4}} = \frac{13}{2} = 6.5$ একক।
    
    2. ধরি, জ্যা $AB$ এর মধ্যবিন্দু $M$। যেহেতু $AB = 12$, তাই $AM = 6$।
    
    \vspace{0.5em}
    3. সমকোণী ত্রিভুজ $\triangle AMC$ হতে কেন্দ্র হতে জ্যা-এর দূরত্ব $CM$ নির্ণয় করি:
    \[ CM = \sqrt{AC^2 - AM^2} = \sqrt{(6.5)^2 - 6^2} = \sqrt{42.25 - 36} = \sqrt{6.25} = 2.5 \text{ একক} \]
    
    4. সমকোণী ত্রিভুজ $\triangle PAC$ (যেখানে $\angle PAC = 90^{\circ}$) এবং জ্যা এর ওপর অঙ্কিত লম্ব $AM$ এর জন্য জ্যামিতিক উপপাদ্য প্রয়োগ করে পাই:
    \[ AC^2 = CM \times PC \implies R^2 = CM \times (CM + PM) \]
    \[ \implies (6.5)^2 = 2.5 \times (2.5 + PM) \implies 42.25 = 2.5 \times (2.5 + PM) \]
    \[ \implies 2.5 + PM = 16.9 \implies PM = 14.4 \text{ একক (এটিই জ্যা হতে $P$ এর দূরত্ব)} \]
    
    5. দূরত্বের ৫ গুণ: 
    \[ 5 \times PM = 5 \times 14.4 = 72 \]
    অতএব, সঠিক উত্তর C।
\end{tcolorbox}

% ------------------------------------------
% QUESTION SECTION
% ------------------------------------------
\begin{tcolorbox}[questionbox={প্রশ্ন (Question) 47}]
    \noindent
    $x^2 + y^2 - 10x - 6y + 30 = 0$ বৃত্তে একটি বর্গক্ষেত্র অন্তর্লিখিত আছে। এই বর্গক্ষেত্রের একটি বাহু $y = x + 3$ সরলরেখার সমান্তরাল। যদি $(x_i, y_i)$ (যেখানে $i=1,2,3,4$) বর্গক্ষেত্রটির শীর্ষবিন্দুসমূহ নির্দেশ করে, তবে $\sum_{i=1}^4 (x_i^2 + y_i^2)$ এর মান কত?
    
    \vspace{0.8em}
    \begin{english}
    \noindent\textit{\color{darkgray}
    A square is inscribed in the circle $x^2 + y^2 - 10x - 6y + 30 = 0$. One side of this square is parallel to the line $y = x + 3$. If $(x_i, y_i)$ for $i=1,2,3,4$ are the vertices of the square, then the value of $\sum_{i=1}^4 (x_i^2 + y_i^2)$ is:
    }
    \end{english}
\end{tcolorbox}

% ------------------------------------------
% OPTIONS SECTION
% ------------------------------------------
\begin{itemize}[label={}, leftmargin=1cm, itemsep=0.5em]
    \item[\textbf{A)}] $148$
    \item[\textbf{B)}] $156$
    \item[\textbf{C)}] $152$
    \item[\textbf{D)}] $160$
\end{itemize}

% ------------------------------------------
% ANSWER HIGHLIGHT
% ------------------------------------------
\vspace{0.5em}
\begin{tcolorbox}[answerbox]
    \textbf{\color{success} উত্তর:} C) $152$
\end{tcolorbox}
\vspace{1em}

% ------------------------------------------
% SOLUTION SECTION
% ------------------------------------------
\begin{tcolorbox}[solutionbox={সমাধান (Solution)}]
    \noindent
    \textbf{ধাপ ১: বৃত্তের সমীকরণটিকে আদর্শ আকারে রূপান্তর} \\
    প্রদত্ত বৃত্তের সমীকরণ: $x^2 + y^2 - 10x - 6y + 30 = 0$
    \[ \implies (x - 5)^2 + (y - 3)^2 = 5^2 + 3^2 - 30 = 4 \]
    সুতরাং, বৃত্তের কেন্দ্র $C(5, 3)$ এবং ব্যাসার্ধ $R = 2$।
    
    \vspace{0.5em}
    \noindent\textbf{ধাপ ২: দূরত্ব সূত্রের জ্যামিতিক বিশ্লেষণ} \\
    ধরি $O(0,0)$ হলো মূলবিন্দু। বৃত্তের পরিধির ওপরিস্থিত বর্গক্ষেত্রের যেকোনো শীর্ষবিন্দু $P_i(x_i, y_i)$ এর জন্য অবস্থান ভেক্টর:
    \[ \vec{OP_i} = \vec{OC} + \vec{CP_i} \]
    উভয়পক্ষকে বর্গ করে পাই:
    \[ OP_i^2 = OC^2 + CP_i^2 + 2\vec{OC} \cdot \vec{CP_i} \]
    
    \vspace{0.5em}
    \noindent\textbf{ধাপ ৩: চারটি শীর্ষবিন্দুর দূরত্বের বর্গের সমষ্টি নির্ণয়} \\
    যেহেতু $P_i$ বৃত্তের ওপর অবস্থিত, তাই $CP_i^2 = R^2 = 4$। চারটি বিন্দুর জন্য সমষ্টি নিলে পাই:
    \[ \sum_{i=1}^4 (x_i^2 + y_i^2) = \sum_{i=1}^4 OP_i^2 = \sum_{i=1}^4 (OC^2 + R^2) + 2\vec{OC} \cdot \left(\sum_{i=1}^4 \vec{CP_i}\right) \]
    
    \vspace{0.5em}
    \noindent\textbf{ধাপ ৪: সাদৃশ্যতার শর্ত প্রয়োগ} \\
    যেহেতু $C$ বিন্দুটি অন্তর্লিখিত বর্গক্ষেত্রের জ্যামিতিক কেন্দ্র (Centroid), তাই কেন্দ্র হতে শীর্ষবিন্দুসমূহের ভেক্টর সমষ্টি শূন্য, অর্থাৎ $\sum_{i=1}^4 \vec{CP_i} = \vec{0}$।
    
    \vspace{0.5em}
    \noindent\textbf{ধাপ ৫: চূড়ান্ত মান হিসাব} \\
    \[ \sum_{i=1}^4 (x_i^2 + y_i^2) = 4OC^2 + 4R^2 \]
    এখানে, $OC^2 = 5^2 + 3^2 = 34$ এবং $R^2 = 4$।
    \[ \therefore \sum_{i=1}^4 (x_i^2 + y_i^2) = 4(34) + 4(4) = 136 + 16 = 152 \]
    অতএব, সঠিক উত্তর C।
\end{tcolorbox}

% ------------------------------------------
% QUESTION SECTION
% ------------------------------------------
\begin{tcolorbox}[questionbox={প্রশ্ন (Question) 48}]
    \noindent
    $x^2 + (y - 1)^2 = 1$ বৃত্তের মূলবিন্দু $O(0,0)$ গামী জ্যা-সমূহের মধ্যবিন্দুসমূহের সঞ্চারপথ সরলরেখা $x + y = 1$ কে যথাক্রমে $P$ ও $Q$ বিন্দুতে ছেদ করে। $PQ$ রেখাংশের দৈর্ঘ্য কত একক?
    
    \vspace{0.8em}
    \begin{english}
    \noindent\textit{\color{darkgray}
    Let the locus of the midpoints of the chords of the circle $x^2 + (y - 1)^2 = 1$ drawn from the origin $O(0,0)$ intersect the line $x + y = 1$ at the points $P$ and $Q$. Then, the length of the line segment $PQ$ is:
    }
    \end{english}
\end{tcolorbox}

% ------------------------------------------
% OPTIONS SECTION
% ------------------------------------------
\begin{itemize}[label={}, leftmargin=1cm, itemsep=0.5em]
    \item[\textbf{A)}] $1$
    \item[\textbf{B)}] $1/\sqrt{2}$
    \item[\textbf{C)}] $\sqrt{2}$
    \item[\textbf{D)}] $1/2$
\end{itemize}

% ------------------------------------------
% ANSWER HIGHLIGHT
% ------------------------------------------
\vspace{0.5em}
\begin{tcolorbox}[answerbox]
    \textbf{\color{success} উত্তর:} B) $1/\sqrt{2}$
\end{tcolorbox}
\vspace{1em}

% ------------------------------------------
% SOLUTION SECTION
% ------------------------------------------
\begin{tcolorbox}[solutionbox={সমাধান (Solution)}]
    \noindent
    \textbf{ধাপ ১: জ্যা-এর মধ্যবিন্দুর সঞ্চারপথ নির্ণয়} \\
    ধরি, মূলবিন্দুগামী যেকোনো জ্যা-এর মধ্যবিন্দু $M(h, k)$। বৃত্তের কেন্দ্র $C(0, 1)$। 
    জ্যামিতিক নিয়ম অনুযায়ী, কেন্দ্র ও জ্যা-এর মধ্যবিন্দুর সংযোজক রেখা জ্যা-এর ওপর লম্ব হয় (অর্থাৎ $CM \perp OM$)। 
    $OM$ এর ঢাল $= \frac{k}{h}$ এবং $CM$ এর ঢাল $= \frac{k - 1}{h}$। যেহেতু তারা লম্ব:
    \[ \left(\frac{k}{h}\right) \times \left(\frac{k - 1}{h}\right) = -1 \implies k(k-1) = -h^2 \implies h^2 + k^2 - k = 0 \]
    চলক প্রতিস্থাপন করে সঞ্চারপথের সমীকরণ পাই: 
    \[ x^2 + y^2 - y = 0 \]
    এটি একটি বৃত্ত যার কেন্দ্র $C'\left(0, \frac{1}{2}\right)$ এবং ব্যাসার্ধ $r = \frac{1}{2}$।
    
    \vspace{0.5em}
    \noindent\textbf{ধাপ ২: রেখা ও সঞ্চারপথের ছেদবিন্দু বিশ্লেষণ} \\
    সঞ্চারপথ বৃত্তটি এবং সরলরেখা $x + y - 1 = 0$ ছেদ করে। নতুন কেন্দ্র $C'\left(0, \frac{1}{2}\right)$ হতে এই সরলরেখাটির লম্ব দূরত্ব $d$:
    \[ d = \frac{\left| 0 + \frac{1}{2} - 1 \right|}{\sqrt{1^2 + 1^2}} = \frac{1}{2\sqrt{2}} \]
    
    \vspace{0.5em}
    \noindent\textbf{ধাপ ৩: জ্যা-এর দৈর্ঘ্য $PQ$ নির্ণয়} \\
    \[ PQ = 2\sqrt{r^2 - d^2} = 2\sqrt{\left(\frac{1}{2}\right)^2 - \left(\frac{1}{2\sqrt{2}}\right)^2} = 2\sqrt{\frac{1}{4} - \frac{1}{8}} = 2\sqrt{\frac{1}{8}} = \frac{1}{\sqrt{2}} \text{ একক} \]
    অতএব, সঠিক উত্তর B।
\end{tcolorbox}

% ------------------------------------------
% QUESTION SECTION
% ------------------------------------------
\begin{tcolorbox}[questionbox={প্রশ্ন (Question) 49}]
    \noindent
    ধরি সরলরেখা $l: \sqrt{2}x + y = a$ প্রথম চতুর্ভাগে অবস্থিত $x^2 + y^2 = 3$ বৃত্ত এবং $x^2 = 2y$ পরাবৃত্তের ছেদবিন্দু $P$ দিয়ে অতিক্রম করে। সরলরেখা $l$ দুটি সমান $2\sqrt{3}$ ব্যাসার্ধের বৃত্ত $C_1$ ও $C_2$ কে স্পর্শ করে। যদি বৃত্তদ্বয়ের কেন্দ্র $Q_1$ ও $Q_2$ যথাক্রমে $y$-অক্ষের ওপর অবস্থিত হয়, তবে ত্রিভুজ $\triangle PQ_1Q_2$ এর ক্ষেত্রফলের বর্গ কত বর্গ একক?
    
    \vspace{0.8em}
    \begin{english}
    \noindent\textit{\color{darkgray}
    Let the line $l: \sqrt{2}x + y = a$ pass through the point of intersection $P$ (in the first quadrant) of the circle $x^2 + y^2 = 3$ and the parabola $x^2 = 2y$. Let the line $l$ touch two circles $C_1$ and $C_2$ of equal radius $2\sqrt{3}$. If the centres $Q_1$ and $Q_2$ of the circles lie on the y-axis, then the square of the area of the triangle $PQ_1Q_2$ is equal to:
    }
    \end{english}
\end{tcolorbox}

% ------------------------------------------
% OPTIONS SECTION
% ------------------------------------------
\begin{itemize}[label={}, leftmargin=1cm, itemsep=0.5em]
    \item[\textbf{A)}] $18$
    \item[\textbf{B)}] $36$
    \item[\textbf{C)}] $72$
    \item[\textbf{D)}] $54$
\end{itemize}

% ------------------------------------------
% ANSWER HIGHLIGHT
% ------------------------------------------
\vspace{0.5em}
\begin{tcolorbox}[answerbox]
    \textbf{\color{success} উত্তর:} C) $72$
\end{tcolorbox}
\vspace{1em}

% ------------------------------------------
% SOLUTION SECTION
% ------------------------------------------
\begin{tcolorbox}[solutionbox={সমাধান (Solution)}]
    \noindent
    \textbf{ধাপ ১: ছেদবিন্দু $P$ নির্ণয়} \\
    পরাবৃত্তের মান বৃত্তের সমীকরণে বসিয়ে পাই:
    \[ 2y + y^2 = 3 \implies y^2 + 2y - 3 = 0 \implies (y + 3)(y - 1) = 0 \]
    যেহেতু বিন্দুটি প্রথম চতুর্ভাগে অবস্থিত, তাই $y = 1$। 
    \[ x^2 = 2(1) = 2 \implies x = \sqrt{2} \]
    ছেদবিন্দুর স্থানাঙ্ক $P(\sqrt{2}, 1)$।
    
    \vspace{0.5em}
    \noindent\textbf{ধাপ ২: সরলরেখা $l$ এর সমীকরণ} \\
    যেহেতু $l$ রেখাটি $P(\sqrt{2}, 1)$ বিন্দুগামী:
    \[ \sqrt{2}(\sqrt{2}) + 1 = a \implies a = 3 \]
    রেখার সমীকরণ: $\sqrt{2}x + y - 3 = 0$।
    
    \vspace{0.5em}
    \noindent\textbf{ধাপ ৩: বৃত্তদ্বয়ের কেন্দ্রদ্বয় নির্ণয়} \\
    ধরি, কেন্দ্রদ্বয় যথাক্রমে $Q_1(0, y_1)$ এবং $Q_2(0, y_2)$। 
    স্পর্শকের শর্তানুযায়ী কেন্দ্র হতে রেখার লম্ব দূরত্ব ব্যাসার্ধ $2\sqrt{3}$ এর সমান:
    \[ \frac{|y_i - 3|}{\sqrt{(\sqrt{2})^2 + 1^2}} = 2\sqrt{3} \implies \frac{|y_i - 3|}{\sqrt{3}} = 2\sqrt{3} \implies |y_i - 3| = 6 \]
    $\therefore y_1 = 9 \implies Q_1(0, 9)$ এবং $y_2 = -3 \implies Q_2(0, -3)$।
    
    \vspace{0.5em}
    \noindent\textbf{ধাপ ৪: ত্রিভুজ $\triangle PQ_1Q_2$ এর ক্ষেত্রফল} \\
    ভূমি $Q_1Q_2 = |9 - (-3)| = 12$ এবং উচ্চতা $P$ বিন্দুর ভুজ $= \sqrt{2}$।
    \[ \text{Area} = \frac{1}{2} \times 12 \times \sqrt{2} = 6\sqrt{2} \]
    \[ \text{Area}^2 = (6\sqrt{2})^2 = 72 \text{ বর্গ একক} \]
    অতএব, সঠিক উত্তর C।
\end{tcolorbox}
% ------------------------------------------
% QUESTION SECTION
% ------------------------------------------
\begin{tcolorbox}[questionbox={প্রশ্ন (Question) 50}]
    \noindent
    বৃত্তদ্বয় $C_1: x^2 + y^2 = 1$ এবং $C_2: x^2 + y^2 - 6x - 8y + 21 = 0$ এর পরিধির ওপরিস্থিত যেকোনো দুটি বিন্দুর মধ্যবর্তী সর্বনিম্ন দূরত্ব কত একক?
    
    \vspace{0.8em}
    \begin{english}
    \noindent\textit{\color{darkgray}
    What is the shortest distance between any two points on the circles $C_1: x^2 + y^2 = 1$ and $C_2: x^2 + y^2 - 6x - 8y + 21 = 0$?
    }
    \end{english}
\end{tcolorbox}

% ------------------------------------------
% OPTIONS SECTION
% ------------------------------------------
\begin{itemize}[label={}, leftmargin=1cm, itemsep=0.5em]
    \item[\textbf{A)}] $1$ একক
    \item[\textbf{B)}] $2$ একক
    \item[\textbf{C)}] $3$ একক
    \item[\textbf{D)}] $4$ একক
\end{itemize}

% ------------------------------------------
% ANSWER HIGHLIGHT
% ------------------------------------------
\vspace{0.5em}
\begin{tcolorbox}[answerbox]
    \textbf{\color{success} উত্তর:} B) $2$ একক
\end{tcolorbox}
\vspace{1em}

% ------------------------------------------
% SOLUTION SECTION
% ------------------------------------------
\begin{tcolorbox}[solutionbox={সমাধান (Solution)}]
    \noindent
    \textbf{ধাপ ১: প্রথম বৃত্তের উপাত্ত বিশ্লেষণ} \\
    প্রথম বৃত্ত $C_1: x^2 + y^2 = 1$ এর কেন্দ্র $O(0,0)$ এবং ব্যাসার্ধ $r_1 = 1$।
    
    \vspace{0.5em}
    \noindent\textbf{ধাপ ২: দ্বিতীয় বৃত্তের উপাত্ত বিশ্লেষণ} \\
    দ্বিতীয় বৃত্ত $C_2: x^2 + y^2 - 6x - 8y + 21 = 0$
    \[ \implies (x - 3)^2 + (y - 4)^2 = 3^2 + 4^2 - 21 = 9 + 16 - 21 = 4 \]
    সুতরাং, দ্বিতীয় বৃত্তের কেন্দ্র $A(3, 4)$ এবং ব্যাসার্ধ $r_2 = \sqrt{4} = 2$।
    
    \vspace{0.5em}
    \noindent\textbf{ধাপ ৩: কেন্দ্রদ্বয়ের মধ্যবর্তী দূরত্ব নির্ণয়} \\
    কেন্দ্রদ্বয়ের মধ্যবর্তী দূরত্ব $d = OA$:
    \[ d = \sqrt{(3 - 0)^2 + (4 - 0)^2} = \sqrt{3^2 + 4^2} = \sqrt{9 + 16} = \sqrt{25} = 5 \]
    
    \vspace{0.5em}
    \noindent\textbf{ধাপ ৪: বৃত্তদ্বয়ের আপেক্ষিক অবস্থান যাচাই} \\
    যেহেতু কেন্দ্রদ্বয়ের দূরত্ব ব্যাসার্ধদ্বয়ের সমষ্টির চেয়ে বেশি (অর্থাৎ $d = 5 > r_1 + r_2 = 1 + 2 = 3$), তাই বৃত্ত দুটি পরস্পর থেকে সম্পূর্ণ পৃথকভাবে বাইরে অবস্থিত।
    
    \vspace{0.5em}
    \noindent\textbf{ধাপ ৫: পরিধিদ্বয়ের মধ্যবর্তী সর্বনিম্ন দূরত্ব হিসাব} \\
    দুটি বহিঃস্থ বৃত্তের পরিধির ওপরিস্থিত বিন্দুদ্বয়ের মধ্যবর্তী সর্বনিম্ন দূরত্ব:
    \[ d_{\min} = d - (r_1 + r_2) = 5 - (1 + 2) = 5 - 3 = 2 \text{ একক} \]
    অতএব, সঠিক উত্তর B।
\end{tcolorbox}
\end{document}